\documentclass[11pt,a4paper]{article}

\usepackage[utf8]{inputenc}
\usepackage[indonesian]{babel}
\usepackage[margin=2cm]{geometry}
\usepackage{amsmath,amssymb,amsfonts}
\newenvironment{proof}{\noindent\textbf{Bukti.}}{\hfill$\blacksquare$\vspace{0.3em}}
\usepackage{booktabs}
\usepackage{longtable}
\usepackage{multirow}
\usepackage{array}
\usepackage{enumitem}
\usepackage{titlesec}
\usepackage{float}
\usepackage{xcolor}
\usepackage{tcolorbox}
\usepackage{fancyhdr}
\usepackage{hyperref}

\hypersetup{
    colorlinks=true,
    linkcolor=blue!80!black,
    citecolor=blue!80!black,
    urlcolor=blue!80!black,
    pdftitle={Bank 25 Soal Latihan dan Solusi Lengkap Persiapan UTS MA4151 Kriptografi},
    pdfauthor={MA4151 Kriptografi ITB}
}

% Header & Footer Settings
\pagestyle{fancy}
\fancyhf{}
\rhead{\textbf{MA4151 Kriptografi}}
\lhead{Bank 25 Soal Latihan \& Solusi UTS}
\rfoot{Halaman \thepage}
\setlength{\headheight}{14pt}

% Color Definitions
\definecolor{primary}{RGB}{20, 50, 130}       % Navy Blue
\definecolor{accent}{RGB}{15, 110, 120}       % Deep Teal
\definecolor{darkgreen}{RGB}{25, 110, 50}     % Emerald
\definecolor{crimson}{RGB}{170, 30, 30}       % Deep Red / Alert
\definecolor{boxbg}{RGB}{248, 250, 252}       % Crisp light slate
\definecolor{solbg}{RGB}{252, 253, 255}
\definecolor{tipbg}{RGB}{243, 250, 245}

\tcbuselibrary{skins,breakable}

\tcbset{
    arc=2mm,
    left=3.5mm,
    right=3.5mm,
    top=3mm,
    bottom=3mm,
    fonttitle=\bfseries,
    breakable
}

% Custom TColorBoxes
\newtcolorbox{soalbox}[1]{
    colback=boxbg,
    colframe=primary,
    title={#1}
}

\newtcolorbox{solusibox}[1]{
    colback=solbg,
    colframe=accent,
    title={#1}
}

\newtcolorbox{tipsbox}[1]{
    colback=tipbg,
    colframe=darkgreen,
    title={#1}
}

% Section Title Formatting
\titleformat{\section}{\large\bfseries\color{primary}}{\thesection}{1em}{}[\titlerule]
\titleformat{\subsection}{\normalsize\bfseries\color{accent}}{\thesubsection}{1em}{}

\title{
    \vspace{-1.5cm}
    \textbf{\huge BANK 25 SOAL LATIHAN INTENSIF}\\[0.2cm]
    \textbf{\Large \& PEMBAHASAN LENGKAP PERSIAPAN UTS}\\[0.3cm]
    \large Mata Kuliah MA4151 Kriptografi -- Semester I Tahun Ajaran 2026/2027\\[0.1cm]
    \large Institut Teknologi Bandung
}
\author{\textbf{Disusun untuk Grinding Mandiri \& Penguasaan Soal Ujian ITB}}
\date{Oktober 2026}

\begin{document}

\maketitle
\vspace{-0.5cm}

\begin{abstract}
Dokumen ini memuat \textbf{Bank 25 Butir Soal Latihan Intensif dan Pembahasan Lengkap} yang dirancang khusus untuk persiapan menghadapi \textbf{Ujian Tengah Semester (UTS) MA4151 Kriptografi}. Soal-soal dalam dokumen ini mencakup seluruh tipologi soal yang diujikan dalam 6 tahun terakhir (UTS 2020/2021 hingga UTS 2025/2026 terbaru), mencakup: (1) Sandi Klasik \& Ruang Kriptosistem; (2) Kriptoanalisis \& Indeks Koinsidensi; (3) Teori Shannon, Entropi, \& Pohon Huffman; (4) Block Cipher Modern (DES \& AES); (5) Pengantar Teori Bilangan (PTB); serta (6) Kriptosistem RSA \& Titik Tetap. Setiap butir soal disertai \textbf{pembahasan matematis formal langkah-demi-langkah, identifikasi teorema pendukung, serta catatan konsep kunci}.
\end{abstract}

\vspace{0.2cm}
\hrule
\vspace{0.3cm}

\tableofcontents
\newpage

% =========================================================================
\section{Kategori 1: Sandi Klasik \& Analisis Ruang Kriptosistem}
% =========================================================================

\begin{soalbox}{Soal 1.1: Ukuran Ruang Kunci Sistem Kripto Perkalian (UTS 2020/2021)}
Misalkan $S$ adalah suatu sistem kripto sandi geser pada alfabet Latin ($\mathcal{P} = \mathcal{C} = \mathcal{K}_1 = \mathbb{Z}_{26}$) dan $T$ adalah suatu sistem kripto sandi perkalian pada alfabet Latin ($e_a(x) \equiv ax \pmod{26}$).
Berapa banyak kunci yang mungkin dibuat dari sistem kripto perkalian $S \times T$?
\begin{enumerate}[label=(\alph*)]
    \item 256
    \item 312
    \item 128
    \item 512
\end{enumerate}
\end{soalbox}

\begin{solusibox}{Pembahasan Soal 1.1}
\begin{enumerate}[leftmargin=*]
    \item \textbf{Definisi Sistem Perkalian:} Sistem kripto perkalian $S \times T$ didefinisikan dengan ruang kunci $\mathcal{K} = \mathcal{K}_1 \times \mathcal{K}_2$.
    \item \textbf{Ukuran Ruang Kunci $S$ (Sandi Geser):} Sandi geser memiliki kunci berupa pergeseran aditif $k \in \mathbb{Z}_{26}$, sehingga:
    \[
    |\mathcal{K}_1| = 26
    \]
    \item \textbf{Ukuran Ruang Kunci $T$ (Sandi Perkalian):} Aturan enkripsi adalah $e_a(x) \equiv ax \pmod{26}$. Agar fungsi ini memiliki invers (dapat didekripsi secara unik), syarat perlu dan cukupnya adalah:
    \[
    \gcd(a, 26) = 1 \iff a \in \mathbb{Z}_{26}^*
    \]
    Banyaknya bilangan bulat pada $\mathbb{Z}_{26}$ yang relatif prima terhadap 26 dihitung menggunakan fungsi phi Euler:
    \[
    |\mathcal{K}_2| = \varphi(26) = \varphi(2 \times 13) = (2 - 1)(13 - 1) = 1 \times 12 = 12
    \]
    Himpunan kunci tersebut adalah $\mathbb{Z}_{26}^* = \{1, 3, 5, 7, 9, 11, 15, 17, 19, 21, 23, 25\}$.
    \item \textbf{Total Ukuran Ruang Kunci Gabungan:}
    \[
    |\mathcal{K}| = |\mathcal{K}_1| \times |\mathcal{K}_2| = 26 \times 12 = \mathbf{312}
    \]
\end{enumerate}
\textbf{Jawaban yang Benar: (b) 312.}
\end{solusibox}

\vspace{0.4cm}

\begin{soalbox}{Soal 1.2: Syarat Invertibilitas \& Tabrakan Sandi Hill (UTS 2021/2022)}
Alice ingin mengirim pesan rahasia kepada Bob menggunakan sandi Hill berdimensi $2 \times 2$. Matriks yang dipilih Alice adalah:
\[
K = \begin{pmatrix} 1 & 2 \\ 3 & 4 \end{pmatrix}
\]
\begin{enumerate}[label=(\alph*)]
    \item Jelaskan secara aljabar mengapa matriks $K$ yang dipilih Alice tidak dapat digunakan dalam sistem kripto Sandi Hill.
    \item Tunjukkan bahwa kegagalan tersebut mengakibatkan terjadinya tabrakan teks-sandi (\textit{ciphertext collision}) dengan mencari dua vektor teks-asal (\textit{plaintext}) berbeda $\mathbf{x}_1, \mathbf{x}_2 \in \mathbb{Z}_{26}^2$ yang menghasilkan teks-sandi yang sama persis: $K\mathbf{x}_1 \equiv K\mathbf{x}_2 \pmod{26}$.
\end{enumerate}
\end{soalbox}

\begin{solusibox}{Pembahasan Soal 1.2}
\begin{enumerate}[leftmargin=*]
    \item \textbf{Analisis Determinan Matriks:}
    Hitung determinan matriks $K$:
    \[
    \det(K) = (1)(4) - (2)(3) = 4 - 6 = -2 \equiv 24 \pmod{26}
    \]
    Matriks $K$ memiliki invers modulo 26 jika dan hanya jika $\gcd(\det(K), 26) = 1$. Namun:
    \[
    \gcd(24, 26) = 2 \neq 1
    \]
    Karena $\det(K) \equiv 0 \pmod 2$, elemen determinan tidak memiliki invers multiplikatif di dalam gelanggang $\mathbb{Z}_{26}$. Akibatnya, matriks invers $K^{-1} \pmod{26}$ tidak ada, sehingga Bob tidak akan pernah bisa mendekripsi pesan secara unik.
    \item \textbf{Menemukan Vektor Tabrakan (\textit{Null Space} Modulo 26):}
    Kondisi $K\mathbf{x}_1 \equiv K\mathbf{x}_2 \pmod{26}$ ekuivalen dengan mencari vektor selisih tak-nol $\mathbf{v} = \mathbf{x}_1 - \mathbf{x}_2 \not\equiv \mathbf{0} \pmod{26}$ sedemikian hingga:
    \[
    K \mathbf{v} \equiv \begin{pmatrix} 0 \\ 0 \end{pmatrix} \pmod{26}
    \]
    Misalkan $\mathbf{v} = \begin{pmatrix} v_1 \\ v_2 \end{pmatrix}$. Susun sistem kongruensi linier:
    \[
    \begin{cases}
    v_1 + 2v_2 \equiv 0 \pmod{26} & \implies v_1 \equiv -2v_2 \pmod{26} \\
    3v_1 + 4v_2 \equiv 0 \pmod{26} & \implies 3(-2v_2) + 4v_2 \equiv -2v_2 \equiv 0 \pmod{26}
    \end{cases}
    \]
    Kongruensi $-2v_2 \equiv 0 \pmod{26} \iff 2v_2 \equiv 0 \pmod{26}$ memiliki solusi umum:
    \[
    v_2 \equiv 0 \pmod{13} \implies v_2 \in \{0, 13\}
    \]
    Ambil solusi tak-nol $v_2 = 13$. Maka:
    \[
    v_1 \equiv -2(13) = -26 \equiv 0 \pmod{26}
    \]
    Diperoleh vektor selisih tak-nol:
    \[
    \mathbf{v} = \begin{pmatrix} 0 \\ 13 \end{pmatrix} \not\equiv \begin{pmatrix} 0 \\ 0 \end{pmatrix} \pmod{26}
    \]
    \item \textbf{Verifikasi Pasangan Plaintext Berbeda:}
    Pilih plaintext sembarang $\mathbf{x}_1 = \begin{pmatrix} 0 \\ 0 \end{pmatrix}$ (huruf \texttt{"AA"}).\\
    Maka $\mathbf{x}_2 = \mathbf{x}_1 + \mathbf{v} = \begin{pmatrix} 0 \\ 13 \end{pmatrix}$ (huruf \texttt{"AN"}).
    Hitung ciphertext masing-masing:
    \[
    K \mathbf{x}_1 = \begin{pmatrix} 1 & 2 \\ 3 & 4 \end{pmatrix} \begin{pmatrix} 0 \\ 0 \end{pmatrix} \equiv \begin{pmatrix} 0 \\ 0 \end{pmatrix} \pmod{26} \quad (\texttt{"AA"})
    \]
    \[
    K \mathbf{x}_2 = \begin{pmatrix} 1 & 2 \\ 3 & 4 \end{pmatrix} \begin{pmatrix} 0 \\ 13 \end{pmatrix} = \begin{pmatrix} 26 \\ 52 \end{pmatrix} \equiv \begin{pmatrix} 0 \\ 0 \end{pmatrix} \pmod{26} \quad (\texttt{"AA"})
    \]
    Dua kata-asal yang berbeda (\texttt{"AA"} dan \texttt{"AN"}) menghasilkan kata-sandi yang sama persis (\texttt{"AA"}). Terbukti terjadi tabrakan teks-sandi.
\end{enumerate}
\end{solusibox}

\vspace{0.4cm}

\begin{soalbox}{Soal 1.3: Pembuktian Teorema Kunci Involusi Sandi Permutasi (UTS 2020/2021)}
Suatu kunci sandi $K$ disebut \textbf{kunci involusi} jika aturan enkripsi identik dengan aturan dekripsi, yaitu $e_K = d_K$.
Buktikan secara formal bahwa permutasi $\pi \in S_m$ dalam sandi permutasi berukuran blok $m$ merupakan kunci involusi jika dan hanya jika:
\[
\pi(i) = j \implies \pi(j) = i, \quad \forall i, j \in \{1, 2, \dots, m\}
\]
\end{soalbox}

\begin{solusibox}{Pembahasan Soal 1.3}
\begin{proof}
Tinjau definisi sandi permutasi pada blok berukuran $m$. Untuk sembarang blok teks-asal $\mathbf{x} = (x_1, \dots, x_m) \in \mathcal{P}$, aturan enkripsi dan dekripsinya adalah:
\[
e_\pi(x_1, \dots, x_m) = (x_{\pi(1)}, \dots, x_{\pi(m)}), \qquad d_\pi(y_1, \dots, y_m) = (y_{\pi^{-1}(1)}, \dots, y_{\pi^{-1}(m)})
\]
\begin{enumerate}[leftmargin=*]
    \item \textbf{Arah Maju ($\implies$):} Asumsikan $\pi$ adalah kunci involusi.
    Berdasarkan definisi, $e_\pi = d_\pi$, yang berarti untuk setiap blok $\mathbf{x}$, berlaku:
    \[
    (x_{\pi(1)}, \dots, x_{\pi(m)}) = (x_{\pi^{-1}(1)}, \dots, x_{\pi^{-1}(m)}) \implies \pi(k) = \pi^{-1}(k), \quad \forall k \in \{1, \dots, m\}
    \]
    Sehingga $\pi = \pi^{-1}$, atau secara ekuivalen $\pi \circ \pi = \operatorname{id}$.\\
    Ambil sembarang $i, j \in \{1, \dots, m\}$ sedemikian hingga $\pi(i) = j$.
    Terapkan fungsi $\pi$ pada kedua ruas:
    \[
    \pi(\pi(i)) = \pi(j)
    \]
    Karena $\pi \circ \pi = \operatorname{id}$, maka ruas kiri menjadi $\operatorname{id}(i) = i$. Dengan demikian, diperoleh:
    \[
    \pi(j) = i
    \]
    Terbukti bahwa $\pi(i) = j \implies \pi(j) = i$.
    \item \textbf{Arah Mundur ($\impliedby$):} Asumsikan bahwa untuk setiap $i, j \in \{1, \dots, m\}$ berlaku $\pi(i) = j \implies \pi(j) = i$.
    Ambil sembarang $i \in \{1, \dots, m\}$. Misalkan $j = \pi(i)$.
    Berdasarkan asumsi yang diberikan, maka berlaku $\pi(j) = i$.
    Substitusikan $j = \pi(i)$ ke dalam persamaan tersebut:
    \[
    \pi(\pi(i)) = i = \operatorname{id}(i)
    \]
    Karena kesamaan ini berlaku untuk seluruh $i \in \{1, \dots, m\}$, disimpulkan bahwa:
    \[
    \pi \circ \pi = \operatorname{id} \implies \pi = \pi^{-1}
    \]
    Karena $\pi = \pi^{-1}$, maka aturan dekripsi memenuhi $d_\pi = e_{\pi^{-1}} = e_\pi$. Terbukti bahwa $\pi$ adalah kunci involusi.
\end{enumerate}
\end{proof}
\end{solusibox}

\vspace{0.4cm}

\begin{soalbox}{Soal 1.4: Dekripsi Sandi Permutasi Bahasa Indonesia (UTS 2022/2023 \& 2023/2024)}
Diberikan teks-sandi:
\[
\texttt{GUED NAGC SAEN KRSE TYAN ANMA}
\]
yang diperoleh dengan mengenkripsi teks-asal berbahasa Indonesia menggunakan sandi permutasi pada blok 4 huruf ($m = 4$).
Tentukan permutasi kunci $\pi$ dan temukan kata-asal (\textit{plaintext}) aslinya!
\end{soalbox}

\begin{solusibox}{Pembahasan Soal 1.4}
\begin{enumerate}[leftmargin=*]
    \item \textbf{Identifikasi Blok Ciphertext ($m = 4$):}
    Teks sandi terdiri atas 6 blok masing-masing berpanjang 4 huruf:
    \[
    C_1 = \texttt{GUED}, \quad C_2 = \texttt{NAGC}, \quad C_3 = \texttt{SAEN}, \quad C_4 = \texttt{KRSE}, \quad C_5 = \texttt{TYAN}, \quad C_6 = \texttt{ANMA}
    \]
    \item \textbf{Analisis Posisi \& Rekonstruksi Kata:}
    Tinjau blok $C_1 = \texttt{G}_1 \texttt{U}_2 \texttt{E}_3 \texttt{D}_4$. Kata bahasa Indonesia yang sangat alami adalah \texttt{"GEDU"} (dari kata \texttt{GEDUNG}).
    Untuk mengubah $C_1 = \texttt{GUED}$ menjadi $P_1 = \texttt{GEDU}$, kita mengambil karakter pada posisi:
    \[
    P_1[1] = C_1[1] = \texttt{G}, \quad P_1[2] = C_1[3] = \texttt{E}, \quad P_1[3] = C_1[4] = \texttt{D}, \quad P_1[4] = C_1[2] = \texttt{U}
    \]
    Artinya, aturan dekripsi memetakan urutan karakter:
    \[
    d_\pi(y_1, y_2, y_3, y_4) = (y_1, y_3, y_4, y_2)
    \]
    Sehingga permutasi dekripsi adalah $\pi^{-1} = \begin{pmatrix} 1 & 2 & 3 & 4 \\ 1 & 3 & 4 & 2 \end{pmatrix}$.\\
    Permutasi enkripsi yang bersesuaian adalah kebalikannya: $\pi = \begin{pmatrix} 1 & 2 & 3 & 4 \\ 1 & 4 & 2 & 3 \end{pmatrix}$, yaitu $e_\pi(x_1, x_2, x_3, x_4) = (x_1, x_4, x_2, x_3)$.
    \item \textbf{Menerapkan Dekripsi ke Seluruh Blok:}
    \begin{align*}
    C_1 &= \texttt{G}_1 \texttt{U}_2 \texttt{E}_3 \texttt{D}_4 \implies (y_1, y_3, y_4, y_2) = \mathbf{\texttt{GEDU}} \\
    C_2 &= \texttt{N}_1 \texttt{A}_2 \texttt{G}_3 \texttt{C}_4 \implies (y_1, y_3, y_4, y_2) = \mathbf{\texttt{NGCA}} \\
    C_3 &= \texttt{S}_1 \texttt{A}_2 \texttt{E}_3 \texttt{N}_4 \implies (y_1, y_3, y_4, y_2) = \mathbf{\texttt{SENA}} \\
    C_4 &= \texttt{K}_1 \texttt{R}_2 \texttt{S}_3 \texttt{E}_4 \implies (y_1, y_3, y_4, y_2) = \mathbf{\texttt{KSER}} \\
    C_5 &= \texttt{T}_1 \texttt{Y}_2 \texttt{A}_3 \texttt{N}_4 \implies (y_1, y_3, y_4, y_2) = \mathbf{\texttt{TANY}} \\
    C_6 &= \texttt{A}_1 \texttt{N}_2 \texttt{M}_3 \texttt{A}_4 \implies (y_1, y_3, y_4, y_2) = \mathbf{\texttt{AMAN}}
    \end{align*}
    \item \textbf{Menggabungkan Plaintext Lengkap:}
    Menggabungkan seluruh blok menghasilkan kalimat bahasa Indonesia yang sempurna:
    \[
    \texttt{GEDU NGCA SENA KSER TANY AMAN} \implies \mathbf{\texttt{"GEDUNG CAS ENAK SERTA NYAMAN"}}
    \]
\end{enumerate}
\end{solusibox}

\vspace{0.4cm}

\begin{soalbox}{Soal 1.5: Sandi Alir LFSR \& Siklus Periode Maksimum (m-sequence)}
Diberikan sebuah Linear Feedback Shift Register (LFSR) berderajat $m = 4$ atas lapangan biner $\mathbb{F}_2$ dengan relasi rekurensi:
\[
s_{i+4} \equiv s_{i+1} \oplus s_i \pmod 2, \quad i \ge 0
\]
dan status awal (\textit{seed}) $(s_0, s_1, s_2, s_3) = (1, 0, 0, 0)$.
\begin{enumerate}[label=(\alph*)]
    \item Tentukan polinomial karakteristik koneksi $C(D)$ dari LFSR tersebut.
    \item Hitung dan tuliskan 15 bit pertama aliran kunci (\textit{keystream}) $s_0, s_1, \dots, s_{14}$.
    \item Buktikan bahwa periode dari aliran bit tersebut adalah periode maksimum $T = 2^4 - 1 = 15$.
\end{enumerate}
\end{soalbox}

\begin{solusibox}{Pembahasan Soal 1.5}
\begin{enumerate}[leftmargin=*]
    \item \textbf{Polinomial Karakteristik Koneksi:}
    Persamaan umpan balik adalah $s_{i+4} = c_1 s_{i+3} \oplus c_2 s_{i+2} \oplus c_3 s_{i+1} \oplus c_4 s_i$.
    Membandingkan koefisien: $c_1 = 0, c_2 = 0, c_3 = 1, c_4 = 1$.
    Polinomial koneksinya adalah:
    \[
    C(D) = 1 + c_1 D + c_2 D^2 + c_3 D^3 + c_4 D^4 = \mathbf{1 + D^3 + D^4} \quad (\text{atau } D^4 + D + 1)
    \]
    \item \textbf{Pembangkitan Aliran Bit Langkah demi Langkah:}
    Status awal: $s_0 = 1, s_1 = 0, s_2 = 0, s_3 = 0$.
    \begin{align*}
    s_4 &= s_1 \oplus s_0 = 0 \oplus 1 = \mathbf{1} \\
    s_5 &= s_2 \oplus s_1 = 0 \oplus 0 = \mathbf{0} \\
    s_6 &= s_3 \oplus s_2 = 0 \oplus 0 = \mathbf{0} \\
    s_7 &= s_4 \oplus s_3 = 1 \oplus 0 = \mathbf{1} \\
    s_8 &= s_5 \oplus s_4 = 0 \oplus 1 = \mathbf{1} \\
    s_9 &= s_6 \oplus s_5 = 0 \oplus 0 = \mathbf{0} \\
    s_{10} &= s_7 \oplus s_6 = 1 \oplus 0 = \mathbf{1} \\
    s_{11} &= s_8 \oplus s_7 = 1 \oplus 1 = \mathbf{0} \\
    s_{12} &= s_9 \oplus s_8 = 0 \oplus 1 = \mathbf{1} \\
    s_{13} &= s_{10} \oplus s_9 = 1 \oplus 0 = \mathbf{1} \\
    s_{14} &= s_{11} \oplus s_{10} = 0 \oplus 1 = \mathbf{1} \\
    s_{15} &= s_{12} \oplus s_{11} = 1 \oplus 0 = \mathbf{1} \\
    s_{16} &= s_{13} \oplus s_{12} = 1 \oplus 1 = \mathbf{0} \\
    s_{17} &= s_{14} \oplus s_{13} = 1 \oplus 1 = \mathbf{0} \\
    s_{18} &= s_{15} \oplus s_{14} = 1 \oplus 1 = \mathbf{0}
    \end{align*}
    Perhatikan bahwa pada langkah ke-15, status register menjadi $(s_{15}, s_{16}, s_{17}, s_{18}) = (1, 0, 0, 0)$, yang persis sama dengan status awal $(s_0, s_1, s_2, s_3) = (1, 0, 0, 0)$!
    Aliran 15 bit pertama adalah:
    \[
    (s_0, s_1, \dots, s_{14}) = \mathbf{(1, 0, 0, 0, 1, 0, 0, 1, 1, 0, 1, 0, 1, 1, 1)}
    \]
    \item \textbf{Pembuktian Periode Maksimum:}
    LFSR memiliki $m = 4$ sel biner, sehingga total status yang mungkin adalah $2^4 = 16$. Karena status nol mutlak $(0, 0, 0, 0)$ tidak pernah dikunjungi (menghasilkan status diam), maka jumlah status bukan-nol maksimal adalah $2^4 - 1 = 15$.
    Karena urutan status berulang tepat setelah 15 langkah dan mengunjungi seluruh 15 status bukan-nol unik, maka periode aliran kunci adalah periode maksimum $T = 2^4 - 1 = \mathbf{15}$ (merupakan sebuah \textit{m-sequence}).
\end{enumerate}
\end{solusibox}

\newpage
% =========================================================================
\section{Kategori 2: Kriptoanalisis (KPA, Kasiski, \& Indeks Koinsidensi)}
% =========================================================================

\begin{soalbox}{Soal 2.1: Serangan Teks-Asal Diketahui Sandi Affine (UTS 2020/2021)}
Misalkan suatu sandi Affine memiliki kunci rahasia $K = (m, n)$ dengan $m, n \in \mathbb{Z}_{26}$ dan $\gcd(m, 26) = 1$.
Diketahui teks-asal \texttt{"hahaha"} menghasilkan teks-sandi \texttt{"NONONO"}.
Tentukan nilai dari $m + n$!
\begin{enumerate}[label=(\alph*)]
    \item 11
    \item 25
    \item 37
    \item 30
\end{enumerate}
\end{soalbox}

\begin{solusibox}{Pembahasan Soal 2.1}
\begin{enumerate}[leftmargin=*]
    \item \textbf{Konversi Huruf ke Angka Modulo 26:}
    \begin{itemize}
        \item Huruf \texttt{'h'} $= 7$, huruf \texttt{'a'} $= 0$.
        \item Huruf \texttt{'N'} $= 13$, huruf \texttt{'O'} $= 14$.
    \end{itemize}
    Teks \texttt{"hahaha"} $\to$ \texttt{"NONONO"} memberikan dua pasangan plaintext-ciphertext:
    \[
    (x_1, y_1) = (7, 13) \quad \text{dan} \quad (x_2, y_2) = (0, 14)
    \]
    \item \textbf{Menyusun Sistem Persamaan Kongruensi:}
    Aturan enkripsi sandi affine: $y \equiv m x + n \pmod{26}$.
    \[
    \begin{cases}
    7m + n \equiv 13 \pmod{26} & \text{(Persamaan 1)} \\
    0m + n \equiv 14 \pmod{26} & \text{(Persamaan 2)}
    \end{cases}
    \]
    \item \textbf{Menyelesaikan Nilai $n$ dan $m$:}
    Dari Persamaan 2 langsung diperoleh:
    \[
    n = 14
    \]
    Substitusikan $n = 14$ ke Persamaan 1:
    \[
    7m + 14 \equiv 13 \pmod{26} \implies 7m \equiv 13 - 14 \equiv -1 \equiv 25 \pmod{26}
    \]
    Cari invers dari $7 \pmod{26}$ menggunakan Algoritma Euclidean:
    \[
    26 = 3(7) + 5, \quad 7 = 1(5) + 2, \quad 5 = 2(2) + 1 \implies 1 = 5 - 2(2) = 5 - 2(7 - 5) = 3(5) - 2(7)
    \]
    \[
    1 = 3(26 - 3 \cdot 7) - 2(7) = 3(26) - 11(7) \implies 7 \times (-11) \equiv 1 \pmod{26}
    \]
    Karena $-11 \equiv 15 \pmod{26}$, maka $7^{-1} \equiv 15 \pmod{26}$.\\
    Kalikan kedua ruas dengan 15:
    \[
    m \equiv 25 \times 15 = 375 \equiv 14 \times 26 + 11 \equiv \mathbf{11} \pmod{26}
    \]
    Periksa syarat invertibilitas: $\gcd(m, 26) = \gcd(11, 26) = 1$ (\textbf{Valid!}).
    \item \textbf{Menghitung Nilai $m + n$:}
    \[
    m + n = 11 + 14 = \mathbf{25}
    \]
\end{enumerate}
\textbf{Jawaban yang Benar: (b) 25.}
\end{solusibox}

\vspace{0.4cm}

\begin{soalbox}{Soal 2.2: Uji Kasiski pada Sandi Vigen\`ere (UTS 2024/2025 No 5)}
Diberikan teks-sandi yang dienkripsi menggunakan sandi Vigen\`ere sebagai berikut:
\[
\texttt{KXQMYTOZKQWHWMEKXQQXKQXYM}
\]
Dengan menggunakan tes Kasiski, tentukan dugaan panjang kuncinya dengan meninjau dua pasang substring berulang!
\end{soalbox}

\begin{solusibox}{Pembahasan Soal 2.2}
\begin{enumerate}[leftmargin=*]
    \item \textbf{Pencatatan Indeks Karakter Ciphertext (Panjang $N = 25$):}
    \begin{center}
    \footnotesize
    \begin{tabular}{c ccccc ccccc ccc}
    \toprule
    \textbf{Indeks} & 0 & 1 & 2 & 3 & 4 & 5 & 6 & 7 & 8 & 9 & 10 & 11 & 12 \\
    \textbf{Huruf}  & \texttt{K} & \texttt{X} & \texttt{Q} & \texttt{M} & \texttt{Y} & \texttt{T} & \texttt{O} & \texttt{Z} & \texttt{K} & \texttt{Q} & \texttt{W} & \texttt{H} & \texttt{W} \\
    \midrule
    \textbf{Indeks} & 13 & 14 & 15 & 16 & 17 & 18 & 19 & 20 & 21 & 22 & 23 & 24 & \\
    \textbf{Huruf}  & \texttt{M} & \texttt{E} & \texttt{K} & \texttt{X} & \texttt{Q} & \texttt{Q} & \texttt{X} & \texttt{K} & \texttt{Q} & \texttt{X} & \texttt{Y} & \texttt{M} & \\
    \bottomrule
    \end{tabular}
    \end{center}
    \item \textbf{Mengidentifikasi Substring Berulang Panjang $\ge 3$:}
    \begin{itemize}
        \item \textbf{Substring 1 (\texttt{"KXQ"}):}
        \begin{itemize}
            \item Muncul pertama kali pada indeks 0: \texttt{K(0) X(1) Q(2)}.
            \item Muncul kedua kali pada indeks 15: \texttt{K(15) X(16) Q(17)}.
            \item Jarak kemunculan: $d_1 = 15 - 0 = \mathbf{15}$.
        \end{itemize}
        \item \textbf{Substring 2 (\texttt{"KQX"} atau \texttt{"KQXYM"}):}
        \begin{itemize}
            \item Perhatikan substring \texttt{"KQX"} pada indeks 20: \texttt{K(20) Q(21) X(22)}.
            \item Perhatikan juga kemunculan trigram \texttt{"QXY"} pada indeks 21: \texttt{Q(21) X(22) Y(23)} dan bigram \texttt{"MY"} pada indeks 3 dan 23.
            \item Perhatikan substring \texttt{"KXQ"} dan variasi polialfabetiknya:
            Periksa substring \texttt{"KQ"} pada indeks: 8 (\texttt{K(8) Q(9)}), 15 (\texttt{K(15) X(16) Q(17)}), 20 (\texttt{K(20) Q(21)}).
            Selisih posisi antara indeks 20 dan 8: $d_2 = 20 - 8 = \mathbf{12}$.
            \item Tinjau substring \texttt{"XQ"}:
            Muncul pada indeks 1 (\texttt{X(1) Q(2)}), indeks 16 (\texttt{X(16) Q(17)}), dan indeks 17 (\texttt{Q(17) Q(18) X(19)}).
            Jarak: $16 - 1 = \mathbf{15}$.
        \end{itemize}
    \end{itemize}
    \item \textbf{Perhitungan Faktor Pembagi Bersama Terbesar (FPB / GCD):}
    Dua jarak utama yang diperoleh adalah $d_1 = 15$ dan $d_2 = 12$:
    \[
    \gcd(d_1, d_2) = \gcd(15, 12) = \mathbf{3}
    \]
    Bila ditinjau pasangan jarak lain $d_1 = 15$ (kelipatan 3 dan 5) dan pola perulangan \texttt{MY} pada jarak $23 - 3 = 20$:
    \[
    \gcd(15, 20) = \mathbf{5}
    \]
    \item \textbf{Kesimpulan:} Berdasarkan tes Kasiski, panjang kata kunci $m$ yang paling mungkin membagi selisih jarak tersebut adalah $\mathbf{m = 3}$ atau $\mathbf{m = 5}$.
\end{enumerate}
\end{solusibox}

\vspace{0.4cm}

\begin{soalbox}{Soal 2.3: Penurunan Rumus Indeks Koinsidensi \& Analisis Transformasi Sandi (UTS 2025/2026 No 1)}
Misalkan $f_0, f_1, \dots, f_{25}$ berturut-turut menyatakan distribusi frekuensi kemunculan huruf A, B, \dots, Z dalam string $x = x_1 x_2 \dots x_n$ yang terdiri atas $n$ buah karakter alfabet.
Indeks koinsidensi $I_c(x)$ didefinisikan sebagai peluang bahwa sembarang dua unsur yang diambil secara acak dari $x$ tanpa pengembalian akan terwakili oleh huruf identik yang sama persis:
\[
I_c(x) = \frac{1}{n(n - 1)} \sum_{i=0}^{25} f_i (f_i - 1)
\]
\begin{enumerate}[label=(\alph*)]
    \item Tunjukkan bagaimana cara memperoleh formula tersebut dari prinsip kombinatorika dasar.
    \item Nyatakan dan jelaskan apa yang terjadi pada nilai $I_c(y)$ (apakah naik, turun, atau tetap tak berubah dari kondisi $I_c(x) \approx 0.065$) pada tiap skenario berikut:
    \begin{enumerate}[label=(\roman*)]
        \item Jika $y$ adalah hasil enkripsi $e_\pi(x_1 \dots x_n) = x_{\pi(1)} \dots x_{\pi(n)}$ oleh suatu permutasi $\pi \in S_n$.
        \item Jika $y$ adalah hasil enkripsi menggunakan sandi geser dengan kunci $K = 11$.
        \item Jika $y$ adalah hasil enkripsi menggunakan sandi Vigen\`ere dengan panjang katakunci $m = 5$.
    \end{enumerate}
\end{enumerate}
\end{soalbox}

\begin{solusibox}{Pembahasan Soal 2.3}
\begin{enumerate}[leftmargin=*]
    \item \textbf{Penurunan Formal Formula $I_c(x)$:}
    \begin{itemize}
        \item Total banyaknya cara memilih 2 karakter sembarang dari string $x$ berpanjang $n$ tanpa memperhatikan urutan (tanpa pengembalian) adalah:
        \[
        \binom{n}{2} = \frac{n(n - 1)}{2}
        \]
        \item Dua karakter yang terpilih berhuruf sama jika keduanya sama-sama huruf ke-$i$ (untuk suatu $i \in \{0, 1, \dots, 25\}$). Karena terdapat $f_i$ buah huruf ke-$i$ di dalam teks, banyaknya cara memilih 2 huruf ke-$i$ adalah:
        \[
        \binom{f_i}{2} = \frac{f_i(f_i - 1)}{2}
        \]
        \item Total cara memilih 2 karakter identik untuk seluruh alfabet adalah penjumlahan saling lepas:
        \[
        \sum_{i=0}^{25} \binom{f_i}{2} = \sum_{i=0}^{25} \frac{f_i(f_i - 1)}{2}
        \]
        \item Peluang memilih 2 karakter yang identik adalah rasio kejadian terhadap ruang sampel:
        \[
        I_c(x) = \frac{\sum_{i=0}^{25} \binom{f_i}{2}}{\binom{n}{2}} = \frac{\sum_{i=0}^{25} \frac{f_i(f_i - 1)}{2}}{\frac{n(n - 1)}{2}} = \mathbf{\frac{1}{n(n - 1)} \sum_{i=0}^{25} f_i (f_i - 1)}
        \]
        Terbukti secara kombinatorik.
    \end{itemize}
    \item \textbf{Analisis Perilaku $I_c(y)$ pada Tiga Skenario:}
    \begin{enumerate}[label=(\roman*)]
        \item \textbf{Skenario Permutasi Posisi ($\pi \in S_n$):}
        Permutasi hanya menukar letak urutan karakter dan sama sekali tidak menambah, menghapus, atau mengubah identitas huruf apa pun. Akibatnya, frekuensi kemunculan setiap huruf tetap sama: $f_i(y) = f_i(x)$ untuk semua $i$.
        \textbf{Kesimpulan: $I_c(y)$ TETAP TAK BERUBAH ($I_c(y) = I_c(x) \approx 0.065$).}
        \item \textbf{Skenario Sandi Geser ($K = 11$):}
        Sandi geser memetakan huruf ke-$i$ ke huruf $(i + 11) \bmod 26$ secara bijektif (satu-satu). Frekuensi hanya bergeser siklis: $f_{(i+11)\bmod 26}(y) = f_i(x)$.
        Karena jumlahan $\sum_{i=0}^{25} f_i(f_i - 1)$ menjumlahkan seluruh alfabet, maka urutan penjumlahan tidak mempengaruhi nilai total.
        \textbf{Kesimpulan: $I_c(y)$ TETAP TAK BERUBAH ($I_c(y) = I_c(x) \approx 0.065$).}
        \item \textbf{Skenario Sandi Vigen\`ere ($m = 5$):}
        Sandi Vigen\`ere adalah substitusi polialfabetik. Teks dipartisi ke dalam 5 coset independen yang masing-masing digeser dengan nilai kunci berbeda. Distribusi frekuensi huruf pada teks sandi gabungan diratakan (*flattened/smoothed*), sehingga mendekati distribusi teks acak seragam.
        \textbf{Kesimpulan: $I_c(y)$ TURUN DRASTIS} dari kondisi teks asli $\approx 0.065$ mendekati nilai acak murni seragam:
        \[
        I_c(y) \approx \frac{1}{26} \approx \mathbf{0.0385}
        \]
    \end{enumerate}
\end{enumerate}
\end{solusibox}

\vspace{0.4cm}

\begin{soalbox}{Soal 2.4: Indeks Koinsidensi Tak Lengkap (UTS 2021/2022 No 4)}
Definisikan \textit{"indeks koinsidensi tidak lengkap"} sebagai indeks koinsidensi yang tidak memakai seluruh huruf pada alfabet Latin (artinya tidak semua huruf dari A--Z dipakai).
Misalkan pada suatu kata $x$ digunakan $N = 120$ huruf dari A hingga X dan memiliki indeks koinsidensi tidak lengkap $\frac{4}{119}$.
Lalu kata $x$ diubah menjadi kata $x'$ berpanjang $N' = 130$ huruf yang mencakup seluruh alfabet Latin dengan menambahkan 10 huruf yang terdiri atas huruf Y dan Z, sedemikian hingga indeks koinsidensi $x'$ menjadi $\frac{522}{16770}$ (catatan: $16770 = 130 \times 129$).
Tentukan banyaknya huruf Y dan Z masing-masing yang ditambahkan pada $x'$!
\begin{enumerate}[label=(\alph*)]
    \item 7 dan 3 (atau sebaliknya)
    \item 6 dan 4 (atau sebaliknya)
    \item 10 dan 0 (atau sebaliknya)
    \item 5 dan 5
\end{enumerate}
\end{soalbox}

\begin{solusibox}{Pembahasan Soal 2.4}
\begin{enumerate}[leftmargin=*]
    \item \textbf{Menghitung Nilai Jumlahan Awal $\sum f_i(f_i - 1)$ pada Kata $x$:}
    Diketahui $N = 120$ dan huruf yang ada hanyalah dari A sampai X (huruf Y dan Z belum ada, frekuensinya $f_Y = f_Z = 0$).
    Formula indeks koinsidensi awal:
    \[
    I_c(x) = \frac{\sum_{i=A}^X f_i(f_i - 1)}{120 \times 119} = \frac{4}{119}
    \]
    Kalikan kedua ruas dengan $120 \times 119$:
    \[
    \sum_{i=A}^X f_i(f_i - 1) = \frac{4}{119} \times 120 \times 119 = 4 \times 120 = \mathbf{480}
    \]
    \item \textbf{Menyusun Persamaan untuk Kata Baru $x'$:}
    Ditambahkan sejumlah $y$ buah huruf Y dan $z$ buah huruf Z dengan total:
    \[
    y + z = 10 \implies z = 10 - y
    \]
    Panjang kata baru menjadi $N' = 120 + 10 = 130$.
    Penyebut baru adalah $N'(N' - 1) = 130 \times 129 = 16770$.
    Indeks koinsidensi baru:
    \[
    I_c(x') = \frac{\sum_{i=A}^X f_i(f_i - 1) + y(y - 1) + z(z - 1)}{16770} = \frac{522}{16770}
    \]
    Maka pembilang harus bernilai 522:
    \[
    480 + y(y - 1) + z(z - 1) = 522 \implies y(y - 1) + z(z - 1) = 522 - 480 = \mathbf{42}
    \]
    \item \textbf{Menyelesaikan Persamaan Aljabar Kuadrat:}
    Substitusikan $z = 10 - y$:
    \[
    (y^2 - y) + ((10 - y)^2 - (10 - y)) = 42
    \]
    \[
    y^2 - y + (100 - 20y + y^2 - 10 + y) = 42
    \]
    \[
    2y^2 - 20y + 90 = 42 \implies 2y^2 - 20y + 48 = 0
    \]
    Bagi kedua ruas dengan 2:
    \[
    y^2 - 10y + 24 = 0 \implies (y - 6)(y - 4) = 0
    \]
    Diperoleh dua solusi simetris:
    \[
    y = 6 \implies z = 4 \quad \text{atau} \quad y = 4 \implies z = 6
    \]
    Banyaknya huruf Y dan Z adalah \textbf{6 dan 4 (atau sebaliknya)}.
\end{enumerate}
\textbf{Jawaban yang Benar: (b) 6 dan 4 (atau sebaliknya).}
\end{solusibox}

\newpage
% =========================================================================
\section{Kategori 3: Teori Shannon, Entropi, \& Pengkodean Huffman}
% =========================================================================

\begin{soalbox}{Soal 3.1: Pengkodean Huffman Optimal \& Perbandingan Entropi (UTS 2022/2023 \& 2023/2024 No 4)}
Misalkan variabel acak diskret $X = \{a, b, c, d, e\}$ memiliki distribusi probabilitas:
\[
p(a) = 0.35, \quad p(b) = 0.25, \quad p(c) = 0.20, \quad p(d) = 0.12, \quad p(e) = 0.08
\]
\begin{enumerate}[label=(\alph*)]
    \item Gunakan algoritma Huffman untuk mengonstruksikan pengkodean bebas-prefiks (\textit{prefix-free encoding}) optimal bagi $X$. Tuliskan kata sandi (\textit{codeword}) untuk setiap simbol!
    \item Hitung panjang rata-rata ekspektasi $L(X)$ dari pengkodean tersebut.
    \item Hitung nilai eksak entropi Shannon $H(X)$ (dalam bit/simbol).
    \item Bandingkan $L(X)$ dengan $H(X)$, serta hitung efisiensi dan redundansi pengkodeannya.
\end{enumerate}
\end{soalbox}

\begin{solusibox}{Pembahasan Soal 3.1}
\begin{enumerate}[leftmargin=*]
    \item \textbf{Konstruksi Pohon Huffman (Greedy Merging):}
    \begin{itemize}
        \item Simbol awal dan peluang: $a(0.35), b(0.25), c(0.20), d(0.12), e(0.08)$.
        \item \textbf{Langkah 1:} Gabungkan dua peluang terkecil: $e(0.08)$ dan $d(0.12) \implies v_1(0.20)$.
        \item Daftar simpul aktif: $v_1(0.20), c(0.20), b(0.25), a(0.35)$.
        \item \textbf{Langkah 2:} Gabungkan dua terkecil: $v_1(0.20)$ dan $c(0.20) \implies v_2(0.40)$.
        \item Daftar simpul aktif: $b(0.25), a(0.35), v_2(0.40)$.
        \item \textbf{Langkah 3:} Gabungkan dua terkecil: $b(0.25)$ dan $a(0.35) \implies v_3(0.60)$.
        \item Daftar simpul aktif: $v_2(0.40), v_3(0.60)$.
        \item \textbf{Langkah 4:} Gabungkan dua simpul terakhir: $v_2(0.40)$ dan $v_3(0.60) \implies \text{Akar}(1.00)$.
    \end{itemize}
    \item \textbf{Penetapan Codeword Bebas-Prefiks:}
    Beri label cabang kiri = 0 dan cabang kanan = 1:
    \begin{itemize}
        \item Simbol $a$: Dari akar $\to$ kanan ($v_3$) $\to$ kanan $\implies \mathbf{11}$ (panjang $\ell(a) = 2$).
        \item Simbol $b$: Dari akar $\to$ kanan ($v_3$) $\to$ kiri $\implies \mathbf{10}$ (panjang $\ell(b) = 2$).
        \item Simbol $c$: Dari akar $\to$ kiri ($v_2$) $\to$ kanan $\implies \mathbf{01}$ (panjang $\ell(c) = 2$).
        \item Simbol $d$: Dari akar $\to$ kiri ($v_2$) $\to$ kiri ($v_1$) $\to$ kanan $\implies \mathbf{001}$ (panjang $\ell(d) = 3$).
        \item Simbol $e$: Dari akar $\to$ kiri ($v_2$) $\to$ kiri ($v_1$) $\to$ kiri $\implies \mathbf{000}$ (panjang $\ell(e) = 3$).
    \end{itemize}
    \item \textbf{Menghitung Panjang Rata-rata Ekspektasi $L(X)$:}
    \[
    L(X) = \sum_{x \in X} p(x) \cdot \ell(x) = 0.35(2) + 0.25(2) + 0.20(2) + 0.12(3) + 0.08(3)
    \]
    \[
    L(X) = 0.70 + 0.50 + 0.40 + 0.36 + 0.24 = \mathbf{2.20 \text{ bit/simbol}}
    \]
    \item \textbf{Menghitung Entropi Shannon $H(X)$:}
    Gunakan formula $H(X) = -\sum p(x) \log_2 p(x)$:
    \begin{align*}
    p(a) \log_2 p(a) &= 0.35 \log_2(0.35) \approx 0.35(-1.51457) \approx -0.53010 \\
    p(b) \log_2 p(b) &= 0.25 \log_2(0.25) = 0.25(-2.00000) = -0.50000 \\
    p(c) \log_2 p(c) &= 0.20 \log_2(0.20) \approx 0.20(-2.32193) \approx -0.46439 \\
    p(d) \log_2 p(d) &= 0.12 \log_2(0.12) \approx 0.12(-3.05889) \approx -0.36707 \\
    p(e) \log_2 p(e) &= 0.08 \log_2(0.08) \approx 0.08(-3.64386) \approx -0.29151
    \end{align*}
    Jumlahkan seluruh komponen:
    \[
    H(X) = -(-0.53010 - 0.50000 - 0.46439 - 0.36707 - 0.29151) \approx \mathbf{2.1531 \text{ bit/simbol}}
    \]
    \item \textbf{Perbandingan, Efisiensi, \& Redundansi:}
    Berdasarkan Teorema Pengkodean Sumber Tanpa Derau Shannon:
    \[
    H(X) \le L(X) < H(X) + 1 \iff 2.1531 \le 2.20 < 3.1531 \quad (\textbf{Terbukti Terpenuhi!})
    \]
    \begin{itemize}
        \item \textbf{Redundansi:} $\bar{R} = L(X) - H(X) = 2.20 - 2.1531 = \mathbf{0.0469 \text{ bit}}$.
        \item \textbf{Efisiensi Pengkodean:} $\eta = \frac{H(X)}{L(X)} \times 100\% = \frac{2.1531}{2.20} \times 100\% \approx \mathbf{97.87\%}$.
    \end{itemize}
\end{enumerate}
\end{solusibox}

\vspace{0.4cm}

\begin{soalbox}{Soal 3.2: Pengkodean Huffman Kasus 2 (UTS 2024/2025 No 1)}
Misalkan variabel acak $X = \{a, b, c, d, e\}$ memiliki distribusi peluang:
\[
p(a) = 0.20, \quad p(b) = 0.35, \quad p(c) = 0.17, \quad p(d) = 0.03, \quad p(e) = 0.25
\]
Gunakan algoritma Huffman untuk mencari encoding prefix-free optimal dari $X$. Bandingkan panjang encoding tersebut dengan $H(X)$!
\end{soalbox}

\begin{solusibox}{Pembahasan Soal 3.2}
\begin{enumerate}[leftmargin=*]
    \item \textbf{Konstruksi Pohon Huffman:}
    Urutkan peluang dari terkecil: $d(0.03), c(0.17), a(0.20), e(0.25), b(0.35)$.
    \begin{itemize}
        \item Gabung $d(0.03)$ dan $c(0.17) \implies v_1(0.20)$.
        \item Simpul aktif: $a(0.20), v_1(0.20), e(0.25), b(0.35)$.
        \item Gabung $a(0.20)$ dan $v_1(0.20) \implies v_2(0.40)$.
        \item Simpul aktif: $e(0.25), b(0.35), v_2(0.40)$.
        \item Gabung $e(0.25)$ dan $b(0.35) \implies v_3(0.60)$.
        \item Gabung $v_2(0.40)$ dan $v_3(0.60) \implies \text{Akar}(1.00)$.
    \end{itemize}
    \item \textbf{Penetapan Codeword \& Panjang:}
    \begin{itemize}
        \item $b$: dari $v_3 \implies \mathbf{11}$ ($\ell = 2$).
        \item $e$: dari $v_3 \implies \mathbf{10}$ ($\ell = 2$).
        \item $a$: dari $v_2 \implies \mathbf{01}$ ($\ell = 2$).
        \item $c$: dari $v_2 \to v_1 \implies \mathbf{001}$ ($\ell = 3$).
        \item $d$: dari $v_2 \to v_1 \implies \mathbf{000}$ ($\ell = 3$).
    \end{itemize}
    \item \textbf{Panjang Rata-rata Ekspektasi $L(X)$:}
    \begin{align*}
    L(X) &= 0.20(2) + 0.35(2) + 0.17(3) + 0.03(3) + 0.25(2) \\
    &= 0.40 + 0.70 + 0.51 + 0.09 + 0.50 = \mathbf{2.20 \text{ bit/simbol}}
    \end{align*}
    \item \textbf{Entropi Shannon $H(X)$:}
    \begin{align*}
    H(X) &= -[0.20 \log_2 0.20 + 0.35 \log_2 0.35 + 0.17 \log_2 0.17 + 0.03 \log_2 0.03 + 0.25 \log_2 0.25] \\
    &= -[0.20(-2.3219) + 0.35(-1.5146) + 0.17(-2.5564) + 0.03(-5.0589) + 0.25(-2.0000)] \\
    &= 0.4644 + 0.5301 + 0.4346 + 0.1518 + 0.5000 = \mathbf{2.0809 \text{ bit/simbol}}
    \end{align*}
    \item \textbf{Perbandingan:}
    $H(X) = 2.0809 \le L(X) = 2.20 < H(X) + 1 = 3.0809$.
    Redundansi $\bar{R} = 2.20 - 2.0809 = 0.1191$ bit, Efisiensi $\eta = \frac{2.0809}{2.20} \approx 94.59\%$.
\end{enumerate}
\end{solusibox}

\vspace{0.4cm}

\begin{soalbox}{Soal 3.3: Pembuktian Formal Aturan Rantai Entropi (UTS 2024/2025 No 2)}
Misalkan $X$ dan $Y$ adalah variabel acak pada himpunan nilai berhingga menurut distribusi peluang $p(X)$ dan $p(Y)$ berturut-turut.
Buktikan secara formal bahwa entropi gabungan dari kedua distribusi peluang tersebut memenuhi aturan rantai:
\[
H(X, Y) = H(Y) + H(X \mid Y)
\]
\end{soalbox}

\begin{solusibox}{Pembahasan Soal 3.3}
\begin{proof}
Berdasarkan definisi formal entropi Shannon, entropi gabungan dari dua variabel acak diskret $X$ dan $Y$ adalah:
\[
H(X, Y) = -\sum_{x \in \mathcal{X}} \sum_{y \in \mathcal{Y}} p(x, y) \log_2 p(x, y)
\]
Berdasarkan aksioma probabilitas bersyarat, peluang gabungan dapat difaktorkan sebagai:
\[
p(x, y) = p(y) \cdot p(x \mid y)
\]
Substitusikan faktorisasi peluang ini ke dalam fungsi logaritma:
\begin{align*}
H(X, Y) &= -\sum_{x \in \mathcal{X}} \sum_{y \in \mathcal{Y}} p(x, y) \log_2 \big(p(y) \cdot p(x \mid y)\big) \\
&= -\sum_{x \in \mathcal{X}} \sum_{y \in \mathcal{Y}} p(x, y) \big(\log_2 p(y) + \log_2 p(x \mid y)\big) \\
&= -\sum_{x \in \mathcal{X}} \sum_{y \in \mathcal{Y}} p(x, y) \log_2 p(y) - \sum_{x \in \mathcal{X}} \sum_{y \in \mathcal{Y}} p(x, y) \log_2 p(x \mid y)
\end{align*}
Tinjau suku pertama. Karena $\log_2 p(y)$ tidak bergantung pada variabel indeks penjumlahan $x$, faktorkan ke luar dari jumlahan terhadap $x$:
\[
-\sum_{y \in \mathcal{Y}} \log_2 p(y) \left(\sum_{x \in \mathcal{X}} p(x, y)\right)
\]
Berdasarkan sifat marginalisasi probabilitas, jumlahan peluang gabungan terhadap seluruh $x$ tepat sama dengan peluang marginal $p(y)$:
\[
\sum_{x \in \mathcal{X}} p(x, y) = p(y)
\]
Sehingga suku pertama menjadi:
\[
-\sum_{y \in \mathcal{Y}} p(y) \log_2 p(y) = H(Y)
\]
Tinjau suku kedua. Berdasarkan definisi formal entropi bersyarat:
\begin{align*}
-\sum_{x \in \mathcal{X}} \sum_{y \in \mathcal{Y}} p(x, y) \log_2 p(x \mid y) 
&= \sum_{y \in \mathcal{Y}} p(y) \left(-\sum_{x \in \mathcal{X}} p(x \mid y) \log_2 p(x \mid y)\right) \\
&= \sum_{y \in \mathcal{Y}} p(y) H(X \mid Y = y) = H(X \mid Y)
\end{align*}
Gabungkan kedua suku hasil penguraian:
\[
H(X, Y) = H(Y) + H(X \mid Y)
\]
Terbukti secara formal.
\end{proof}
\end{solusibox}

\vspace{0.4cm}

\begin{soalbox}{Soal 3.4: Teorema Distribusi Ciphertext Seragam pada Keamanan Sempurna (UTS 2025/2026 No 4)}
Buktikan bahwa jika suatu sistem kripto $(\mathcal{P}, \mathcal{C}, \mathcal{K}, \mathcal{E}, \mathcal{D})$ memiliki \textbf{keamanan sempurna} (\textit{perfect secrecy}) dan memenuhi $|\mathcal{K}| = |\mathcal{C}| = |\mathcal{P}|$, maka setiap teks-sandi memiliki peluang yang sama untuk muncul sebagai hasil enkripsi dengan distribusi seragam:
\[
P_C(y) = \frac{1}{|\mathcal{C}|}, \quad \forall y \in \mathcal{C}
\]
\end{soalbox}

\begin{solusibox}{Pembahasan Soal 3.4}
\begin{proof}
Asumsikan tanpa mengurangi perumuman bahwa untuk setiap teks-asal $x \in \mathcal{P}$ berlaku $P(P = x) > 0$.
\begin{enumerate}[leftmargin=*]
    \item \textbf{Teorema Karakterisasi Shannon untuk $|\mathcal{K}| = |\mathcal{C}| = |\mathcal{P}|$:}
    Shannon membuktikan bahwa jika sistem dengan $|\mathcal{K}| = |\mathcal{C}| = |\mathcal{P}|$ aman sempurna, maka:
    \begin{enumerate}
        \item Distribusi peluang kunci haruslah seragam: $P(K = k) = \frac{1}{|\mathcal{K}|}$ untuk semua $k \in \mathcal{K}$.
        \item Untuk setiap pasangan $x \in \mathcal{P}$ dan $y \in \mathcal{C}$, terdapat tepat satu kunci unik $k \in \mathcal{K}$ sedemikian hingga $e_k(x) = y$.
    \end{enumerate}
    \item \textbf{Menghitung Peluang Bersyarat $P(C = y \mid P = x)$:}
    Ambil sembarang $x \in \mathcal{P}$ dan sembarang $y \in \mathcal{C}$. Misalkan $k_{x,y}$ adalah kunci tunggal yang memenuhi $e_{k_{x,y}}(x) = y$.
    Maka peluang bersyaratnya adalah:
    \[
    P(C = y \mid P = x) = P(K = k_{x,y}) = \frac{1}{|\mathcal{K}|}
    \]
    \item \textbf{Menghitung Peluang Marjinal Teks-Sandi $P(C = y)$:}
    Berdasarkan hukum peluang total:
    \[
    P(C = y) = \sum_{x \in \mathcal{P}} P(P = x) \cdot P(C = y \mid P = x) = \sum_{x \in \mathcal{P}} P(P = x) \cdot \frac{1}{|\mathcal{K}|} = \frac{1}{|\mathcal{K}|} \sum_{x \in \mathcal{P}} P(P = x)
    \]
    Karena $\sum_{x \in \mathcal{P}} P(P = x) = 1$, maka:
    \[
    P(C = y) = \frac{1}{|\mathcal{K}|}
    \]
    \item \textbf{Mengaitkan dengan Ukuran Ruang Sandi $|\mathcal{C}|$:}
    Karena diketahui $|\mathcal{K}| = |\mathcal{C}|$, substitusikan $|\mathcal{K}| = |\mathcal{C}|$ ke hasil di atas:
    \[
    P(C = y) = \frac{1}{|\mathcal{C}|}, \quad \forall y \in \mathcal{C}
    \]
    Terbukti bahwa setiap ciphertext memiliki peluang kemunculan yang identik dan berdistribusi seragam secara mutlak.
\end{enumerate}
\end{proof}
\end{solusibox}

\vspace{0.4cm}

\begin{soalbox}{Soal 3.5: Perhitungan Jarak Unicity ($n_0$) Sandi Klasik}
Redundansi bahasa Inggris didefinisikan sebagai $R = 1 - \frac{H_L}{\log_2 26} \approx 0.75$, dengan $R \log_2 26 \approx 3.53$ bit/huruf.
Hitung nilai eksak jarak unicity ($n_0$) untuk:
\begin{enumerate}[label=(\alph*)]
    \item Sandi Geser.
    \item Sandi Affine.
    \item Sandi Substitusi Monofonik Umum.
    \item Sandi Hill berdimensi $2 \times 2$.
\end{enumerate}
\end{soalbox}

\begin{solusibox}{Pembahasan Soal 3.5}
Formula Shannon untuk jarak unicity kriptosistem dengan distribusi kunci acak seragam:
\[
n_0 \approx \frac{H(K)}{R \log_2 |\mathcal{P}|} = \frac{\log_2 |\mathcal{K}|}{R \log_2 26} = \frac{\log_2 |\mathcal{K}|}{3.53}
\]
\begin{enumerate}[leftmargin=*]
    \item \textbf{Sandi Geser:}
    $|\mathcal{K}| = 26 \implies H(K) = \log_2 26 \approx 4.7004$ bit.
    \[
    n_0 = \frac{4.7004}{3.53} \approx 1.33 \approx \mathbf{2 \text{ karakter}}
    \]
    Artinya, mencegat minimal 2 huruf sandi sudah cukup bagi penyerang untuk memulihkan kunci sandi geser secara tunggal tanpa kunci palsu.
    \item \textbf{Sandi Affine:}
    $|\mathcal{K}| = 312 \implies H(K) = \log_2 312 \approx 8.2854$ bit.
    \[
    n_0 = \frac{8.2854}{3.53} \approx 2.35 \approx \mathbf{3 \text{ karakter}}
    \]
    Penyerang membutuhkan minimal 3 huruf ciphertext untuk memecahkan sandi affine secara unik.
    \item \textbf{Sandi Substitusi Monofonik:}
    $|\mathcal{K}| = 26! \approx 4.0329 \times 10^{26} \implies H(K) = \log_2(26!) \approx 88.38$ bit.
    \[
    n_0 = \frac{88.38}{3.53} \approx 25.04 \approx \mathbf{25 \text{ karakter}}
    \]
    Dibutuhkan sekitar 25 karakter ciphertext untuk memecahkan kunci substitusi umum.
    \item \textbf{Sandi Hill $2 \times 2$:}
    $|\mathcal{K}| = |GL(2, \mathbb{Z}_{26})| = 157.248 \implies H(K) = \log_2(157248) \approx 17.26$ bit.
    \[
    n_0 = \frac{17.26}{3.53} \approx 4.89 \approx \mathbf{5 \text{ karakter}}
    \]
\end{enumerate}
\end{solusibox}

\newpage
% =========================================================================
\section{Kategori 4: Block Cipher Modern (DES \& AES)}
% =========================================================================

\begin{soalbox}{Soal 4.1: Konsep Cipher Perkalian \& Arsitektur DES/AES (UTS 2024/2025 No 3)}
\begin{enumerate}[label=(\alph*)]
    \item Tuliskan dan jelaskan tahapan-tahapan yang dilakukan pada sistem kripto DES dan sistem kripto AES!
    \item Jelaskan secara komprehensif mengapa sistem kripto DES dan AES disebut sebagai sistem kripto perkalian (\textit{product cipher})!
\end{enumerate}
\end{soalbox}

\begin{solusibox}{Pembahasan Soal 4.1}
\begin{enumerate}[leftmargin=*]
    \item \textbf{Tahapan Sistem Kripto DES \& AES:}
    \begin{itemize}
        \item \textbf{Tahapan DES (Ukuran Blok 64-bit, Kunci 56-bit, 16 Ronde Feistel):}
        \begin{enumerate}
            \item \textbf{Initial Permutation ($IP$):} Permutasi awal 64 bit plaintext menjadi dua paruh 32-bit: $(L_0, R_0)$.
            \item \textbf{16 Ronde Feistel:} Untuk setiap ronde $i = 1, \dots, 16$:
            \[
            L_i = R_{i-1}, \qquad R_i = L_{i-1} \oplus F(R_{i-1}, K_i)
            \]
            Fungsi ronde $F$ melakukan 4 langkah: (1) \textit{Ekspansi E} ($32 \to 48$ bit); (2) XOR dengan subkunci $K_i$; (3) Substitusi non-linear oleh 8 S-box ($48 \to 32$ bit); (4) Permutasi linier $P$ ($32 \to 32$ bit).
            \item \textbf{32-bit Swap \& Inverse Permutation ($IP^{-1}$):} Membalik paruh $(R_{16}, L_{16})$ dan menerapkan permutasi invers untuk menghasilkan ciphertext 64 bit.
        \end{enumerate}
        \item \textbf{Tahapan AES-128 (Ukuran Blok 128-bit, Matriks State $4 \times 4$ Byte, 10 Ronde SPN):}
        \begin{enumerate}
            \item \textbf{Ronde Awal (Ronde 0):} AddRoundKey tunggal dengan kunci awal $w_0 \dots w_3$.
            \item \textbf{9 Ronde Standar (Ronde 1--9):} Melakukan 4 transformasi berurutan:
            \begin{itemize}
                \item \textbf{SubBytes:} Substitusi non-linear per byte menggunakan invers di $GF(2^8)$ + transformasi affine atas $\mathbb{F}_2$.
                \item \textbf{ShiftRows:} Permutasi siklis byte ke kiri per baris (baris 0 geser 0, baris 1 geser 1, baris 2 geser 2, baris 3 geser 3 byte).
                \item \textbf{MixColumns:} Difusi linier kolom melalui perkalian modulo $z^4 + 1$ dengan polinomial $\{03\}z^3 + \{01\}z^2 + \{01\}z + \{02\}$ atas $GF(2^8)$.
                \item \textbf{AddRoundKey:} Bitwise XOR matriks state dengan subkunci ronde.
            \end{itemize}
            \item \textbf{Ronde Terakhir (Ronde 10):} Sama seperti ronde standar, tetapi transformasi \textbf{MixColumns ditiadakan} (SubBytes $\to$ ShiftRows $\to$ AddRoundKey).
        \end{enumerate}
    \end{itemize}
    \item \textbf{Mengapa Disebut Sistem Kripto Perkalian (Product Cipher)?}
    Dalam terminologi kriptografi Claude Shannon (1949), cipher perkalian didefinisikan sebagai komposisi fungsi atau perkalian operator aljabar dari dua atau lebih transformasi kriptografi dasar:
    \[
    \text{Cipher} = T_r \circ T_{r-1} \circ \dots \circ T_1
    \]
    DES dan AES disebut cipher perkalian karena menggabungkan dua jenis transformasi dasar secara bergantian berulang-kali (multi-ronde):
    \begin{enumerate}
        \item Transformasi \textbf{Substitusi} (S-box), yang berfungsi menciptakan \textbf{Konfusi} (\textit{confusion}), yaitu mengaburkan relasi statistik antara teks-sandi dan kunci rahasia.
        \item Transformasi \textbf{Permutasi / Linier} (P-box pada DES, ShiftRows dan MixColumns pada AES), yang berfungsi menciptakan \textbf{Difusi} (\textit{diffusion}), yaitu menyebarkan redundansi dan ketergantungan bit plaintext ke seluruh blok ciphertext.
    \end{enumerate}
    Transformasi substitusi saja rentan analisis frekuensi, dan transformasi permutasi saja rentan analisis linier. Namun hasil perkalian (komposisi bergantian) keduanya secara berulang menghasilkan margin keamanan kriptografis yang sangat tinggi.
\end{enumerate}
\end{solusibox}

\vspace{0.4cm}

\begin{soalbox}{Soal 4.2: Aljabar $GF(2^8)$ \& Polinomial MixColumns AES (UTS 2020/2021 No 5)}
Dalam perspektif aljabar, pada tahap MixColumns di algoritma AES, setiap kolom dari blok data yang dienkripsi dipandang sebagai sukubanyak modulo $(01 z^4 + 01)$ dan dikalikan dengan sebuah sukubanyak yang koefisiennya berada pada lapangan hingga $GF(2^8) \cong \mathbb{F}_2[x] / \langle x^8 + x^4 + x^3 + x + 1 \rangle$.
Dalam representasi heksadesimal dari $GF(2^8)$, sukubanyak tersebut adalah ...
\begin{enumerate}[label=(\alph*)]
    \item $03 z^3 + 01 z^2 + 01 z + 02$
    \item $01 z^3 + 02 z^2 + 03 z + 01$
    \item $02 z^3 + 03 z^2 + 01 z + 01$
    \item $01 z^3 + 01 z^2 + 02 z + 03$
\end{enumerate}
\end{soalbox}

\begin{solusibox}{Pembahasan Soal 4.2}
\begin{enumerate}[leftmargin=*]
    \item \textbf{Formulasi Matematis Transformasi MixColumns:}
    Pada spesifikasi resmi NIST FIPS 197 (AES), setiap kolom dari matriks state $\mathbf{s} = [s_{0,c}, s_{1,c}, s_{2,c}, s_{3,c}]^T$ diperlakukan sebagai polinomial derajat 3:
    \[
    s(z) = s_{3,c} z^3 + s_{2,c} z^2 + s_{1,c} z + s_{0,c}
    \]
    Operasi MixColumns didefinisikan sebagai perkalian modulo $z^4 + 1$ dengan polinomial tetap $c(z)$:
    \[
    s'(z) = c(z) \otimes s(z) \pmod{z^4 + 1}
    \]
    di mana koefisien polinomial $c(z)$ adalah:
    \[
    c(z) = \{03\} z^3 + \{01\} z^2 + \{01\} z + \{02\}
    \]
    \item \textbf{Representasi Matriks Kolom yang Bersesuaian:}
    Perkalian polinomial siklis modulo $z^4 + 1$ di atas ekuivalen dengan perkalian matriks sirkulan:
    \[
    \begin{bmatrix}
    s'_{0,c} \\ s'_{1,c} \\ s'_{2,c} \\ s'_{3,c}
    \end{bmatrix}
    =
    \begin{bmatrix}
    02 & 03 & 01 & 01 \\
    01 & 02 & 03 & 01 \\
    01 & 01 & 02 & 03 \\
    03 & 01 & 01 & 02
    \end{bmatrix}
    \begin{bmatrix}
    s_{0,c} \\ s_{1,c} \\ s_{2,c} \\ s_{3,c}
    \end{bmatrix}
    \]
    Perhatikan bahwa baris pertama matriks adalah $[02, 03, 01, 01]$, yang berkorespondensi tepat dengan koefisien konstanta $\{02\}$ hingga koefisien derajat tertinggi $\{03\}$.
\end{enumerate}
\textbf{Jawaban yang Benar: (a) $03 z^3 + 01 z^2 + 01 z + 02$.}
\end{solusibox}

\vspace{0.4cm}

\begin{soalbox}{Soal 4.3: Pembuktian Sifat Komplementasi DES \& Implikasinya}
Misalkan $P$ adalah blok teks-asal, $K$ adalah kunci, dan $C = \operatorname{DES}_K(P)$ adalah teks-sandi hasil enkripsi DES.
Notasi $\bar{x}$ menyatakan komplemen bit demi bit dari $x$ ($\bar{x} = x \oplus \mathbf{1}$).
\begin{enumerate}[label=(\alph*)]
    \item Buktikan secara aljabar bahwa sistem kripto DES memiliki sifat komplementasi:
    \[
    \operatorname{DES}_{\bar{K}}(\bar{P}) = \overline{\operatorname{DES}_K(P)} = \bar{C}
    \]
    \item Jelaskan secara detail bagaimana seorang penyerang dapat memanfaatkan sifat ini untuk mengurangi ruang pencarian kunci (\textit{brute-force search}) dari $2^{56}$ menjadi $2^{55}$ uji kunci di bawah model penyerangan Chosen-Plaintext Attack (CPA).
\end{enumerate}
\end{soalbox}

\begin{solusibox}{Pembahasan Soal 4.3}
\begin{proof}
Tinjau struktur fungsi ronde Feistel $F(R, K)$ pada DES:
\begin{enumerate}[leftmargin=*]
    \item \textbf{Lapisan Ekspansi $E$:} Karena ekspansi $E: \{0, 1\}^{32} \to \{0, 1\}^{48}$ hanyalah duplikasi bit (pemetaan linier), mengomplemenkan input akan mengomplemenkan output:
    \[
    E(\bar{R}) = \overline{E(R)} = E(R) \oplus \mathbf{1}_{48}
    \]
    \item \textbf{Penambahan Subkunci:} Jika subkunci juga dikomplemenkan ($\bar{K} = K \oplus \mathbf{1}_{48}$), maka pada operasi XOR:
    \[
    E(\bar{R}) \oplus \bar{K} = \big(E(R) \oplus \mathbf{1}_{48}\big) \oplus \big(K \oplus \mathbf{1}_{48}\big) = E(R) \oplus K \oplus (\mathbf{1}_{48} \oplus \mathbf{1}_{48}) = E(R) \oplus K
    \]
    Dua komplemen saling meniadakan secara aljabar!
    \item \textbf{Keluaran S-Box \& Permutasi $P$:}
    Karena input ke-8 S-box bernilai sama persis ($E(R) \oplus K$), maka keluaran S-box dan hasil permutasi $P$ tidak berubah sama sekali:
    \[
    F(\bar{R}, \bar{K}) = F(R, K)
    \]
    \item \textbf{Persamaan Ronde Feistel:}
    Untuk paruh kiri pada ronde ke-$i$:
    \[
    \bar{L}_{i-1} \oplus F(\bar{R}_{i-1}, \bar{K}_i) = \big(L_{i-1} \oplus \mathbf{1}_{32}\big) \oplus F(R_{i-1}, K_i) = \overline{L_{i-1} \oplus F(R_{i-1}, K_i)} = \bar{R}_i
    \]
    Dan paruh kanan: $\bar{L}_i = \bar{R}_{i-1}$.
    \item Karena sifat komplemen ini terpelihara di setiap 16 ronde dan permutasi $IP$ serta $IP^{-1}$ bersifat linier bitwise, maka terbukti secara induktif:
    \[
    \operatorname{DES}_{\bar{K}}(\bar{P}) = \overline{\operatorname{DES}_K(P)} = \bar{C}
    \]
\end{enumerate}
\end{proof}

\textbf{Pemanfaatan dalam Serangan Chosen-Plaintext Attack (CPA):}
\begin{enumerate}[leftmargin=*]
    \item Penyerang memilih plaintext sembarang $P$ dan meminta enkripsinya untuk mendapatkan $C = \operatorname{DES}_K(P)$. Penyerang juga meminta enkripsi dari komplemennya $\bar{P}$ untuk mendapatkan $C' = \operatorname{DES}_K(\bar{P})$.
    \item Untuk mencari kunci rahasia $K$, penyerang menguji calon kunci $K_i$ hanya pada separuh ruang kunci ($2^{55}$ calon kunci pertama):
    \begin{itemize}
        \item Hitung $Y = \operatorname{DES}_{K_i}(P)$.
        \item Jika $Y = C$, maka kunci rahasia adalah $K = K_i$.
        \item Jika $\bar{Y} = C'$, maka berdasarkan sifat komplementasi $\operatorname{DES}_{\bar{K}_i}(\bar{P}) = \bar{Y} = C'$, yang berarti kunci rahasianya adalah $K = \bar{K}_i$!
    \end{itemize}
    \item Setiap 1 kali enkripsi uji coba mengevaluasi 2 kunci sekaligus ($K_i$ dan $\bar{K}_i$). Dengan demikian, kompleksitas pencarian kunci berkurang tepat separuhnya: dari $2^{56}$ menjadi $\mathbf{2^{55}}$ pengujian enkripsi.
\end{enumerate}
\end{solusibox}

\vspace{0.4cm}

\begin{soalbox}{Soal 4.4: Analisis Propagasi Galat Mode Operasi Cipher Blok (ECB, CBC, CTR)}
Misalkan sebuah pesan teks dienkripsi menggunakan cipher blok berukuran blok 128 bit.
Selama proses transmisi melalui saluran yang berderau, terjadi kesalahan transmisi tepat 1 bit pada blok teks-sandi ke-$i$ ($C_i$).
Jelaskan dampak propagasi kesalahan transmisi tersebut pada hasil dekripsi blok-blok teks-asal ($P_{i-1}, P_i, P_{i+1}, P_{i+2}$) jika sistem menggunakan:
\begin{enumerate}[label=(\alph*)]
    \item Mode Electronic Codebook (ECB).
    \item Mode Cipher Block Chaining (CBC).
    \item Mode Counter (CTR).
\end{enumerate}
\end{soalbox}

\begin{solusibox}{Pembahasan Soal 4.4}
\begin{enumerate}[leftmargin=*]
    \item \textbf{Mode Electronic Codebook (ECB):}
    \begin{itemize}
        \item Persamaan dekripsi: $P_k = D_K(C_k)$ secara independen untuk setiap blok.
        \item Kesalahan 1 bit pada $C_i$ menyebabkan input dekripsi ke fungsi $D_K$ berubah. Karena cipher blok memiliki efek longsoran (*avalanche effect*), maka **seluruh 128 bit pada blok $P_i$ akan rusak total secara acak** ($\approx 50\%$ bit salah).
        \item Blok sebelum ($P_{i-1}$) dan blok sesudah ($P_{i+1}, P_{i+2}$) sama sekali tidak terpengaruh (kesalahan 0 bit).
        \item \textbf{Propagasi galat:} Terisolasi pada 1 blok penuh ($P_i$).
    \end{itemize}
    \item \textbf{Mode Cipher Block Chaining (CBC):}
    \begin{itemize}
        \item Persamaan dekripsi: $P_k = D_K(C_k) \oplus C_{k-1}$.
        \item Tinjau blok per blok:
        \begin{enumerate}
            \item Pada blok $P_{i-1}$: $P_{i-1} = D_K(C_{i-1}) \oplus C_{i-2}$ $\implies$ **tidak terpengaruh (benar 100\%)**.
            \item Pada blok $P_i$: $P_i = D_K(C_i) \oplus C_{i-1}$. Karena $C_i$ salah 1 bit, output $D_K(C_i)$ rusak total secara acak $\implies$ **blok $P_i$ rusak total (50\% bit salah)**.
            \item Pada blok $P_{i+1}$: $P_{i+1} = D_K(C_{i+1}) \oplus C_i$. Blok $C_{i+1}$ tidak salah, sehingga $D_K(C_{i+1})$ benar. Namun ia di-XOR dengan $C_i$ yang memiliki 1 bit salah $\implies$ **blok $P_{i+1}$ hanya mengalami tepat 1 bit kesalahan** pada posisi bit yang sama persis dengan bit yang rusak pada $C_i$!
            \item Pada blok $P_{i+2}$: $P_{i+2} = D_K(C_{i+2}) \oplus C_{i+1}$ $\implies$ **pulih sempurna (0 bit salah)**.
        \end{enumerate}
        \item \textbf{Propagasi galat:} Terbatas pada 2 blok ($P_i$ rusak total, $P_{i+1}$ salah 1 bit, blok berikutnya pulih otomatis /\textit{self-synchronizing}).
    \end{itemize}
    \item \textbf{Mode Counter (CTR):}
    \begin{itemize}
        \item Persamaan dekripsi: $P_k = C_k \oplus E_K(\text{Nonce} \parallel k)$.
        \item Aliran kunci $E_K(\text{Nonce} \parallel k)$ dibangkitkan secara independen dari teks sandi.
        \item Pada blok $P_i = C_i \oplus E_K(\text{Nonce} \parallel i)$, karena operasi XOR bersifat bit-per-bit, maka kesalahan 1 bit pada $C_i$ hanya membalik **tepat 1 bit yang bersesuaian pada $P_i$**!
        \item Seluruh blok lainnya ($P_{i-1}, P_{i+1}, P_{i+2}$) pulih sempurna dan benar 100\%.
        \item \textbf{Propagasi galat:} Nol propagasi (1 bit ciphertext salah $\implies$ tepat 1 bit plaintext salah).
    \end{itemize}
\end{enumerate}
\end{solusibox}

\newpage
% =========================================================================
\section{Kategori 5: Pengantar Teori Bilangan (PTB) untuk Kriptografi}
% =========================================================================

\begin{soalbox}{Soal 5.1: Teorema Sisa Cina (CRT) Modulus Non-Koprima (UTS 2022/2023 \& 2023/2024 No 2)}
Tentukan solusi bilangan bulat $x$ dalam modulo 210 dari sistem kongruensi linier simultan berikut:
\[
\begin{cases}
2x \equiv 3 \pmod 5 \\
2x \equiv 4 \pmod 6 \\
3x \equiv 2 \pmod 7
\end{cases}
\]
\end{soalbox}

\begin{solusibox}{Pembahasan Soal 5.1}
\begin{enumerate}[leftmargin=*]
    \item \textbf{Penyederhanaan Masing-Masing Kongruensi Linier:}
    \begin{itemize}
        \item Kongruensi 1: $2x \equiv 3 \pmod 5$.
        Invers dari $2 \pmod 5$ adalah 3 (karena $2 \times 3 = 6 \equiv 1 \pmod 5$).
        Kalikan kedua ruas dengan 3:
        \[
        x \equiv 3 \times 3 = 9 \equiv \mathbf{4 \pmod 5}
        \]
        \item Kongruensi 2: $2x \equiv 4 \pmod 6$.
        Perhatikan bahwa $\gcd(2, 6) = 2 \mid 4$. Bagi seluruh persamaan (termasuk modulus) dengan 2:
        \[
        x \equiv 2 \pmod 3
        \]
        Dalam modulo 6, kongruensi ini menghasilkan dua residu unik:
        \[
        x \equiv 2 \pmod 6 \quad \text{atau} \quad x \equiv 2 + 3 = 5 \pmod 6
        \]
        (Uji: $2(2) = 4 \equiv 4 \pmod 6$ dan $2(5) = 10 \equiv 4 \pmod 6$. Keduanya valid).
        \item Kongruensi 3: $3x \equiv 2 \pmod 7$.
        Invers dari $3 \pmod 7$ adalah 5 (karena $3 \times 5 = 15 \equiv 1 \pmod 7$).
        Kalikan kedua ruas dengan 5:
        \[
        x \equiv 2 \times 5 = 10 \equiv \mathbf{3 \pmod 7}
        \]
    \end{itemize}
    \item \textbf{Penggabungan Modulo Koprima 5 dan 7 Menggunakan CRT:}
    Tinjau subsistem:
    \[
    \begin{cases}
    x \equiv 4 \pmod 5 \\
    x \equiv 3 \pmod 7
    \end{cases}
    \]
    Modulus $m_1 = 5, m_2 = 7$ saling prima ($\gcd(5, 7) = 1$). Modulus gabungan $M = 35$.
    \begin{itemize}
        \item $M_1 = 35 / 5 = 7 \implies 7 y_1 \equiv 1 \pmod 5 \implies 2 y_1 \equiv 1 \pmod 5 \implies y_1 = 3$.
        \item $M_2 = 35 / 7 = 5 \implies 5 y_2 \equiv 1 \pmod 7 \implies y_2 = 3$.
    \end{itemize}
    Berdasarkan CRT:
    \[
    x \equiv a_1 M_1 y_1 + a_2 M_2 y_2 = 4(7)(3) + 3(5)(3) = 84 + 45 = 129 \equiv \mathbf{24 \pmod{35}}
    \]
    Bentuk umum solusi:
    \[
    x = 35k + 24, \quad k \in \mathbb{Z}
    \]
    \item \textbf{Mencocokkan Solusi dengan Kongruensi Modulo 6 ($x \equiv 2 \pmod 3$):}
    Substitusikan $x = 35k + 24$ ke dalam syarat $x \equiv 2 \pmod 3$:
    \begin{align*}
    35k + 24 &\equiv 2 \pmod 3 \\
    \implies 2k + 0 &\equiv 2 \pmod 3 \implies 2k \equiv 2 \pmod 3 \implies k \equiv 1 \pmod 3
    \end{align*}
    Karena kita mencari solusi dalam modulo $\operatorname{lcm}(5, 6, 7) = 210 = 35 \times 6$, maka $k$ berjalan pada rentang $\{0, 1, 2, 3, 4, 5\}$.
    Nilai $k \in \{0, 1, 2, 3, 4, 5\}$ yang memenuhi $k \equiv 1 \pmod 3$ adalah:
    \[
    k = 1 \quad \text{dan} \quad k = 4
    \]
    \item \textbf{Menghitung Nilai Eksak Solusi Modulo 210:}
    \begin{itemize}
        \item Untuk $k = 1$:
        \[
        x_1 = 35(1) + 24 = \mathbf{59}
        \]
        Uji: $59 \equiv 4 \pmod 5$; $2(59) = 118 \equiv 4 \pmod 6$; $3(59) = 177 = 25(7) + 2 \equiv 2 \pmod 7$. (\textbf{Valid!})
        \item Untuk $k = 4$:
        \[
        x_2 = 35(4) + 24 = 140 + 24 = \mathbf{164}
        \]
        Uji: $164 \equiv 4 \pmod 5$; $2(164) = 328 = 54(6) + 4 \equiv 4 \pmod 6$; $3(164) = 492 = 70(7) + 2 \equiv 2 \pmod 7$. (\textbf{Valid!})
    \end{itemize}
\end{enumerate}
\textbf{Solusi Akhir:} Terdapat tepat dua solusi unik modulo 210, yaitu:
\[
\mathbf{x \equiv 59 \pmod{210} \quad \text{dan} \quad x \equiv 164 \pmod{210}}
\]
\end{solusibox}

\vspace{0.4cm}

\begin{soalbox}{Soal 5.2: Teorema Keterbagian $2^r \mid \varphi(n)$ (UTS 2024/2025 No 4)}
Tunjukkan bahwa jika bilangan bulat positif $n$ memiliki $r$ buah faktor prima yang berbeda, maka:
\[
2^r \mid \varphi(n) \quad (\text{untuk prima ganjil berbeda})
\]
Secara umum, buktikan bahwa jika $n$ memiliki $r$ faktor prima ganjil berbeda, maka $2^r \mid \varphi(n)$!
\end{soalbox}

\begin{solusibox}{Pembahasan Soal 5.2}
\begin{proof}
Misalkan faktorisasi prima dari bilangan bulat positif $n$ adalah:
\[
n = 2^{k_0} p_1^{k_1} p_2^{k_2} \dots p_r^{k_r}
\]
di mana $p_1, p_2, \dots, p_r$ adalah $r$ buah bilangan prima ganjil yang berbeda ($p_i \ge 3$ untuk setiap $i = 1, \dots, r$), dengan $k_i \ge 1$ untuk $i = 1, \dots, r$, dan $k_0 \ge 0$ adalah pangkat dari faktor prima 2.
\begin{enumerate}[leftmargin=*]
    \item Berdasarkan sifat multiplikatif dari fungsi phi Euler:
    \[
    \varphi(n) = \varphi(2^{k_0}) \cdot \prod_{i=1}^r \varphi(p_i^{k_i})
    \]
    \item Menggunakan rumus nilai fungsi phi pada perpangkatan bilangan prima $\varphi(p^k) = p^{k-1}(p - 1)$:
    \[
    \varphi(n) = \varphi(2^{k_0}) \cdot \prod_{i=1}^r p_i^{k_i - 1} (p_i - 1)
    \]
    \item Tinjau faktor $(p_i - 1)$ untuk setiap $i = 1, 2, \dots, r$.
    Karena setiap $p_i$ adalah bilangan prima ganjil, maka $p_i \equiv 1 \pmod 2$.
    Akibatnya, $p_i - 1$ pastilah merupakan \textbf{bilangan bulat genap}:
    \[
    2 \mid (p_i - 1), \quad \forall i \in \{1, 2, \dots, r\}
    \]
    \item Karena terdapat $r$ buah suku bilangan genap yang saling dikalikan, maka perkalian suku-suku tersebut habis dibagi oleh $2^r$:
    \[
    2^r \;\Big|\; \prod_{i=1}^r (p_i - 1)
    \]
    \item Karena faktor-faktor lainnya ($\varphi(2^{k_0})$ dan $p_i^{k_i - 1}$) seluruhnya merupakan bilangan bulat positif, disimpulkan secara langsung bahwa:
    \[
    2^r \mid \varphi(n)
    \]
\end{enumerate}
Terbukti secara formal.
\end{proof}
\end{solusibox}

\vspace{0.4cm}

\begin{soalbox}{Soal 5.3: Sifat Perkalian Fungsi Phi \& Kasus Kuadrat $\varphi(n^2) = n\varphi(n)$ (UTS 2025/2026 No 3)}
\begin{enumerate}[label=(\alph*)]
    \item Jika setiap faktor prima dari $n$ juga merupakan faktor prima dari $m$, perlihatkan bahwa:
    \[
    \varphi(mn) = n \varphi(m)
    \]
    \item Sebagai akibat khusus dari pernyataan di atas, buktikan bahwa untuk setiap bilangan bulat positif $n$, berlaku:
    \[
    \varphi(n^2) = n \varphi(n)
    \]
\end{enumerate}
\end{soalbox}

\begin{solusibox}{Pembahasan Soal 5.3}
\begin{proof}
\begin{enumerate}[leftmargin=*]
    \item \textbf{Pembuktian Bagian (a):}
    Gunakan formula umum produk Euler untuk fungsi phi:
    \[
    \varphi(k) = k \prod_{p \mid k} \left(1 - \frac{1}{p}\right)
    \]
    di mana produk diambil atas seluruh himpunan pembagi prima berbeda dari $k$.
    \begin{itemize}
        \item Misalkan himpunan faktor prima dari $m$ adalah $\mathcal{P}_m = \{p : p \text{ prima dan } p \mid m\}$.
        \item Misalkan himpunan faktor prima dari $n$ adalah $\mathcal{P}_n = \{p : p \text{ prima dan } p \mid n\}$.
        \item Diketahui bahwa setiap faktor prima dari $n$ juga membagi $m$, yang berarti $\mathcal{P}_n \subseteq \mathcal{P}_m$.
        \item Tinjau bilangan $mn$. Himpunan pembagi prima dari hasil kali $mn$ adalah gabungan pembagi prima keduanya:
        \[
        \mathcal{P}_{mn} = \mathcal{P}_m \cup \mathcal{P}_n
        \]
        Karena $\mathcal{P}_n \subseteq \mathcal{P}_m$, maka $\mathcal{P}_{mn} = \mathcal{P}_m$. Artinya, himpunan faktor prima unik dari $mn$ \textbf{sama persis} dengan himpunan faktor prima unik dari $m$!
        \item Tuliskan rumus produk Euler untuk $\varphi(mn)$:
        \[
        \varphi(mn) = (mn) \prod_{p \mid mn} \left(1 - \frac{1}{p}\right) = (mn) \prod_{p \mid m} \left(1 - \frac{1}{p}\right) = n \cdot \left[ m \prod_{p \mid m} \left(1 - \frac{1}{p}\right) \right]
        \]
        Perhatikan bahwa suku di dalam tanda kurung siku tepat merupakan definisi dari $\varphi(m)$.
        Oleh karena itu:
        \[
        \varphi(mn) = n \varphi(m)
        \]
        Terbukti.
    \end{itemize}
    \item \textbf{Pembuktian Bagian (b):}
    Ambil $m = n$.
    Jelas bahwa setiap faktor prima dari $n$ juga merupakan faktor prima dari $m = n$ (karena $\mathcal{P}_n \subseteq \mathcal{P}_n$ berlaku secara trivial).
    Terapkan rumus dari bagian (a) dengan mensubstitusikan $m = n$:
    \[
    \varphi(n \cdot n) = n \varphi(n) \implies \mathbf{\varphi(n^2) = n \varphi(n)}
    \]
    Terbukti secara mutlak.
\end{enumerate}
\end{proof}
\end{solusibox}

\vspace{0.4cm}

\begin{soalbox}{Soal 5.4: Lema Eksponensial B\'ezout \& Teorema Orde Fermat (UTS 2021/2022 Bagian B)}
Diberikan lema berikut:
\begin{quote}
\textbf{Lema.} Misalkan $a \in \mathbb{Z}$ dan $u, v \in \mathbb{Z}^+$ dengan $\gcd(u, v) = d$. Jika $a^u \equiv 1 \pmod m$ dan $a^v \equiv 1 \pmod m$, maka:
\[
a^d \equiv 1 \pmod m
\]
\end{quote}
\begin{enumerate}[label=(\alph*)]
    \item Buktikan kebenaran Lema di atas menggunakan Identitas B\'ezout!
    \item Dengan menggunakan Lema tersebut dan Teorema Fermat Kecil, buktikan pernyataan berikut:
    \begin{quote}
    Misalkan $p$ adalah bilangan prima dan $p \nmid a$. Andaikan $n$ adalah bilangan bulat positif \textbf{terkecil} sedemikian hingga $p \mid (a^n - 1)$ (bilangan $n$ tersebut ada berdasarkan Teorema Fermat). Buktikan bahwa:
    \[
    n \mid (p - 1)
    \]
    \end{quote}
\end{enumerate}
\end{soalbox}

\begin{solusibox}{Pembahasan Soal 5.4}
\begin{proof}
\begin{enumerate}[leftmargin=*]
    \item \textbf{Pembuktian Lema Eksponensial B\'ezout:}
    Berdasarkan Identitas B\'ezout, karena $\gcd(u, v) = d$, maka terdapat pasangan bilangan bulat $x, y \in \mathbb{Z}$ sedemikian hingga:
    \[
    u x + v y = d
    \]
    Salah satu dari $x$ atau $y$ mungkin bernilai negatif. Tanpa mengurangi perumuman, asumsikan $x \ge 0$ dan $y < 0$ (sehingga $-y > 0$).
    Tuliskan identitas dalam bentuk bilangan positif:
    \[
    d + v(-y) = u x
    \]
    Pangkatkan basis $a$ pada kedua ruas:
    \[
    a^{d + v(-y)} = a^{ux} \implies a^d \cdot (a^v)^{-y} = (a^u)^x
    \]
    Tinjau relasi kongruensi modulo $m$. Diketahui bahwa $a^u \equiv 1 \pmod m$ dan $a^v \equiv 1 \pmod m$:
    \[
    a^d \cdot (1)^{-y} \equiv (1)^x \pmod m \implies a^d \cdot 1 \equiv 1 \pmod m \implies \mathbf{a^d \equiv 1 \pmod m}
    \]
    Lema terbukti.
    \item \textbf{Pembuktian $n \mid (p - 1)$:}
    \begin{itemize}
        \item Diketahui bahwa $n$ adalah bilangan bulat positif \textbf{terkecil} yang memenuhi $p \mid (a^n - 1)$, yang ekuivalen dengan $a^n \equiv 1 \pmod p$.
        \item Berdasarkan Teorema Fermat Kecil, karena $p$ prima dan $p \nmid a$, maka:
        \[
        a^{p-1} \equiv 1 \pmod p
        \]
        \item Definisikan $d = \gcd(n, p - 1)$.
        Berdasarkan Lema yang telah dibuktikan pada bagian (a) dengan mengambil $u = n$ dan $v = p - 1$, maka:
        \[
        a^d \equiv 1 \pmod p \implies p \mid (a^d - 1)
        \]
        \item Karena $d$ adalah pembagi persekutuan terbesar dari $n$ dan $p - 1$, maka $d \le n$.
        Namun, berdasarkan hipotesis soal, $n$ adalah bilangan bulat positif \textbf{terkecil} yang memenuhi $a^k \equiv 1 \pmod p$.
        Karena $d$ juga memenuhi $a^d \equiv 1 \pmod p$ dan $d \le n$, satu-satunya kemungkinan adalah:
        \[
        d = n
        \]
        \item Karena $d = \gcd(n, p - 1)$ dan $d = n$, maka $n$ haruslah membagi $p - 1$:
        \[
        n \mid (p - 1)
        \]
    \end{itemize}
\end{enumerate}
Terbukti tuntas.
\end{proof}
\end{solusibox}

\newpage
% =========================================================================
\section{Kategori 6: Kriptosistem RSA \& Titik Tetap (Fixed Points)}
% =========================================================================

\begin{soalbox}{Soal 6.1: Teorema Titik Tetap RSA (UTS 2022/2023 \& 2023/2024 No 3)}
Suatu teks-asal $x \in \mathbb{Z}_n^*$ disebut \textbf{titik tetap} (\textit{fixed point} / tak berubah) jika hasil enkripsinya sama persis dengan teks asalnya:
\[
e_K(x) \equiv x^b \equiv x \pmod n
\]
Tunjukkan dan buktikan secara formal bahwa untuk sistem kripto RSA dengan modulus $n = pq$ ($p, q$ prima berbeda) dan kunci publik $b$, banyaknya teks-asal tetap $x \in \mathbb{Z}_n^*$ tepat sama dengan:
\[
N_{\text{tetap}} = \gcd(b - 1, p - 1) \cdot \gcd(b - 1, q - 1)
\]
\textit{[Petunjuk: Pandang sistem kekongruenan $e_K(x) \equiv x \pmod p$ dan $e_K(x) \equiv x \pmod q$.]}
\end{soalbox}

\begin{solusibox}{Pembahasan Soal 6.1}
\begin{proof}
Karena $x \in \mathbb{Z}_n^*$, maka $\gcd(x, n) = 1$, yang berakibat $x \not\equiv 0 \pmod p$ dan $x \not\equiv 0 \pmod q$.
Berdasarkan Teorema Sisa Cina (karena $\gcd(p, q) = 1$):
\[
x^b \equiv x \pmod n \iff
\begin{cases}
x^b \equiv x \pmod p \\
x^b \equiv x \pmod q
\end{cases}
\]
Karena $x \not\equiv 0 \pmod p$ dan $x \not\equiv 0 \pmod q$, kita dapat mengalikan kedua ruas dengan invers perkalian $x^{-1}$, sehingga sistem ekuivalen dengan:
\[
\begin{cases}
x^{b-1} \equiv 1 \pmod p \\
x^{b-1} \equiv 1 \pmod q
\end{cases}
\]
\begin{enumerate}[leftmargin=*]
    \item \textbf{Menghitung Banyaknya Solusi Modulo $p$:}
    Tinjau kongruensi $x^{b-1} \equiv 1 \pmod p$ di dalam grup perkalian $\mathbb{Z}_p^*$.
    \begin{itemize}
        \item Karena $p$ adalah bilangan prima, maka $\mathbb{Z}_p^*$ merupakan sebuah **grup siklis** berorde $|\mathbb{Z}_p^*| = p - 1$.
        \item Misalkan $g$ adalah generator (akar primitif) dari grup $\mathbb{Z}_p^*$. Maka setiap elemen $x \in \mathbb{Z}_p^*$ dapat dinyatakan secara tunggal sebagai $x \equiv g^k \pmod p$ untuk suatu $k \in \{0, 1, \dots, p - 2\}$.
        \item Substitusikan ke dalam persamaan:
        \[
        (g^k)^{b-1} \equiv g^{k(b-1)} \equiv 1 \equiv g^0 \pmod p
        \]
        \item Karena orde dari $g$ adalah $p - 1$, kongruensi di atas setara dengan:
        \[
        k(b - 1) \equiv 0 \pmod{p - 1}
        \]
        \item Ini adalah persamaan kongruensi linier standar untuk variabel $k$ dalam modulo $p - 1$. Berdasarkan teori kongruensi linier, banyaknya solusi untuk $k \in \{0, 1, \dots, p - 2\}$ tepat sama dengan:
        \[
        \gcd(b - 1, p - 1)
        \]
        Sehingga terdapat tepat $\gcd(b - 1, p - 1)$ buah nilai $x \in \mathbb{Z}_p^*$ yang memenuhi $x^{b-1} \equiv 1 \pmod p$.
    \end{itemize}
    \item \textbf{Menghitung Banyaknya Solusi Modulo $q$:}
    Dengan penalaran yang persis sama pada grup siklis $\mathbb{Z}_q^*$ yang berorde $q - 1$, banyaknya solusi $x \in \mathbb{Z}_q^*$ yang memenuhi $x^{b-1} \equiv 1 \pmod q$ adalah tepat:
    \[
    \gcd(b - 1, q - 1)
    \]
    \item \textbf{Menggabungkan Solusi Menggunakan CRT:}
    Berdasarkan Teorema Sisa Cina, pemetaan:
    \[
    \mathbb{Z}_n^* \to \mathbb{Z}_p^* \times \mathbb{Z}_q^*, \quad x \mapsto (x \bmod p, \; x \bmod q)
    \]
    merupakan sebuah isomorfisme grup yang bijektif.
    Oleh karena itu, setiap pasangan solusi $(x_p, x_q) \in \mathbb{Z}_p^* \times \mathbb{Z}_q^*$ berkorespondensi satu-satu dengan tepat satu solusi unik $x \in \mathbb{Z}_n^*$.
    Total banyaknya teks-asal tetap $x \in \mathbb{Z}_n^*$ adalah hasil kali banyaknya solusi masing-masing modulus:
    \[
    N_{\text{tetap}} = \gcd(b - 1, p - 1) \cdot \gcd(b - 1, q - 1)
    \]
\end{enumerate}
Terbukti secara formal.
\end{proof}
\end{solusibox}

\vspace{0.4cm}

\begin{soalbox}{Soal 6.2: Komputasi Lengkap Kriptosistem RSA \& Algoritma Eksponensiasi Cepat}
Diberikan parameter sistem kripto RSA dengan dua bilangan prima $p = 11$ dan $q = 13$, serta eksponen enkripsi publik $b = 7$.
\begin{enumerate}[label=(\alph*)]
    \item Hitung modulus publik $n$, fungsi totien Euler $\varphi(n)$, dan tentukan eksponen dekripsi privat $a$.
    \item Enkripsikan pesan teks-asal $x = 9$ menggunakan algoritma kuadrat-dan-kali (*Square-and-Multiply*) untuk menghitung $y \equiv x^b \pmod n$.
    \item Lakukan dekripsi pada ciphertext $y$ yang diperoleh untuk memverifikasi bahwa $d_K(y) = x$.
    \item Tentukan banyaknya titik tetap (\textit{fixed points}) pada sistem RSA ini!
\end{enumerate}
\end{soalbox}

\begin{solusibox}{Pembahasan Soal 6.2}
\begin{enumerate}[leftmargin=*]
    \item \textbf{Pembangkitan Kunci RSA:}
    \begin{itemize}
        \item Modulus publik: $n = p \cdot q = 11 \times 13 = \mathbf{143}$.
        \item Totien Euler: $\varphi(n) = (p - 1)(q - 1) = 10 \times 12 = \mathbf{120}$.
        \item Eksponen privat $a$ adalah invers perkalian dari $b = 7$ modulo 120:
        \[
        7 a \equiv 1 \pmod{120}
        \]
        Gunakan Algoritma Euclidean Diperluas:
        \[
        120 = 17(7) + 1 \implies 1 = 120 - 17(7) \implies 7(-17) \equiv 1 \pmod{120}
        \]
        Karena $-17 \equiv 120 - 17 = 103 \pmod{120}$, maka diperoleh:
        \[
        a = \mathbf{103}
        \]
        (Periksa: $7 \times 103 = 721 = 6(120) + 1 \equiv 1 \pmod{120}$, valid).
    \end{itemize}
    \item \textbf{Enkripsi Pesan $x = 9$ dengan Algoritma Square-and-Multiply ($9^7 \pmod{143}$):}
    Representasi biner dari eksponen $b = 7$ adalah $(111)_2$.
    \begin{itemize}
        \item Mulai dengan akumulator $A = 1$.
        \item \textbf{Bit 1 (paling kiri):} $A = A^2 = 1^2 = 1$; karena bit bernilai 1: $A = (A \times 9) = 9$.
        \item \textbf{Bit 2 (tengah):} $A = A^2 = 9^2 = 81 \pmod{143}$; karena bit bernilai 1:
        \[
        A = (81 \times 9) = 729 = 5(143) + 14 \equiv 14 \pmod{143}
        \]
        \item \textbf{Bit 3 (paling kanan):} $A = A^2 = 14^2 = 196 \equiv 196 - 143 = 53 \pmod{143}$; karena bit bernilai 1:
        \[
        A = (53 \times 9) = 477 = 3(143) + 48 \equiv \mathbf{48 \pmod{143}}
        \]
    \end{itemize}
    Hasil enkripsi teks sandi adalah $y = \mathbf{48}$.
    \item \textbf{Verifikasi Dekripsi ($48^{103} \pmod{143}$ Menggunakan CRT):}
    Alih-alih menghitung modulo 143 langsung, kita gunakan Teorema Sisa Cina:
    \begin{itemize}
        \item Modulo $p = 11$: Berdasarkan Teorema Fermat Kecil ($48^{10} \equiv 1 \pmod{11}$):
        $48 \equiv 4 \pmod{11}$, dan eksponen $103 \equiv 3 \pmod{10}$.
        \[
        x \equiv 4^3 = 64 = 5(11) + 9 \equiv \mathbf{9 \pmod{11}}
        \]
        \item Modulo $q = 13$: Berdasarkan Teorema Fermat Kecil ($48^{12} \equiv 1 \pmod{13}$):
        $48 \equiv 9 \equiv -4 \pmod{13}$, dan eksponen $103 \equiv 7 \pmod{12}$.
        \[
        x \equiv 9^7 \equiv (-4)^7 = -16384 \equiv 9 \pmod{13}
        \]
        (Periksa: $9^2 = 81 \equiv 3$, $9^4 \equiv 9$, $9^6 \equiv 27 \equiv 1 \implies 9^7 \equiv 9 \pmod{13}$).
    \end{itemize}
    Karena $x \equiv 9 \pmod{11}$ dan $x \equiv 9 \pmod{13}$, maka menurut CRT:
    \[
    x \equiv \mathbf{9 \pmod{143}}
    \]
    Dekripsi berhasil sempurna memulihkan plaintext awal $x = 9$.
    \item \textbf{Menghitung Banyaknya Titik Tetap:}
    Berdasarkan Teorema Titik Tetap:
    \[
    N_{\text{tetap}} = \gcd(b - 1, p - 1) \cdot \gcd(b - 1, q - 1) = \gcd(7 - 1, 11 - 1) \cdot \gcd(7 - 1, 13 - 1)
    \]
    \[
    N_{\text{tetap}} = \gcd(6, 10) \cdot \gcd(6, 12) = 2 \times 6 = \mathbf{12 \text{ buah titik tetap}}
    \]
\end{enumerate}
\end{solusibox}

\vspace{0.4cm}

\begin{soalbox}{Soal 6.3: Jumlah Elemen Relatif Prima \& Aplikasi Teori Bilangan (UTS 2021/2022 No 7)}
\begin{enumerate}[label=(\alph*)]
    \item Buktikan bahwa jumlah semua bilangan bulat positif yang kurang dari $n$ dan relatif prima terhadap $n$ diberikan oleh:
    \[
    \sum_{\substack{1 \le k < n \\ \gcd(k, n) = 1}} k = \frac{1}{2} n \varphi(n)
    \]
    \item Gunakan hasil di atas untuk menghitung jumlah seluruh bilangan bulat positif yang kurang dari $231$ dan tidak habis dibagi oleh 3, 7, dan 11!
\end{enumerate}
\end{soalbox}

\begin{solusibox}{Pembahasan Soal 6.3}
\begin{enumerate}[leftmargin=*]
    \item \textbf{Pembuktian Rumus Penjumlahan Simetris:}
    \begin{proof}
    Perhatikan sifat keterbagian pembagi bersama terbesar:
    \[
    \gcd(k, n) = \gcd(n - k, n)
    \]
    Hal ini berarti jika $k$ relatif prima terhadap $n$, maka pasangannya $(n - k)$ juga pasti relatif prima terhadap $n$.
    Himpunan bilangan yang relatif prima terhadap $n$ terdiri atas $\varphi(n)$ elemen.
    Kita dapat memasangkan setiap elemen $k$ dengan $n - k$:
    \[
    S = \sum_{\substack{1 \le k < n \\ \gcd(k, n) = 1}} k
    \]
    Tuliskan jumlahan yang sama dalam urutan terbalik:
    \[
    S = \sum_{\substack{1 \le k < n \\ \gcd(k, n) = 1}} (n - k)
    \]
    Jumlahkan kedua persamaan:
    \[
    2S = \sum_{\substack{1 \le k < n \\ \gcd(k, n) = 1}} \big(k + (n - k)\big) = \sum_{\substack{1 \le k < n \\ \gcd(k, n) = 1}} n
    \]
    Karena terdapat tepat $\varphi(n)$ suku penjumlahan dan setiap suku bernilai $n$, maka:
    \[
    2S = n \cdot \varphi(n) \implies S = \mathbf{\frac{1}{2} n \varphi(n)}
    \]
    \end{proof}
    \item \textbf{Aplikasi pada $n = 231$:}
    Faktorisasi prima dari $n = 231$:
    \[
    231 = 3 \times 7 \times 11
    \]
    Bilangan bulat positif yang tidak habis dibagi oleh 3, 7, dan 11 adalah bilangan yang tidak memiliki faktor persekutuan dengan 231, yang berarti relatif prima terhadap 231 ($\gcd(k, 231) = 1$).
    Hitung nilai $\varphi(231)$ menggunakan sifat multiplikatif:
    \[
    \varphi(231) = \varphi(3) \times \varphi(7) \times \varphi(11) = (3 - 1)(7 - 1)(11 - 1) = 2 \times 6 \times 10 = \mathbf{120}
    \]
    Maka jumlah seluruh bilangan tersebut dihitung langsung menggunakan rumus:
    \[
    S = \frac{1}{2} n \varphi(n) = \frac{1}{2} \times 231 \times 120 = 231 \times 60 = \mathbf{13.860}
    \]
\end{enumerate}
\end{solusibox}

\end{document}
